\documentclass[11pt]{article}
\usepackage[margin=1in]{geometry}
\usepackage[T1]{fontenc}
\usepackage{lmodern}
\usepackage{amsmath,amssymb,amsthm,mathtools}
\usepackage{booktabs,array,longtable}
\usepackage{xcolor,listings}
\usepackage{xurl}
\usepackage[hidelinks]{hyperref}
\hypersetup{pdftitle={From an Uninformative Spectral Bound to Useful Stopping Guarantees for QC-GHP Codes},pdfauthor={}}
\newtheorem{theorem}{Theorem}
\newtheorem*{priorresult}{Parameter obstruction}
\setlength{\emergencystretch}{3em}
\lstset{basicstyle=\scriptsize\ttfamily,columns=fullflexible,keepspaces=true,
 breaklines=true,breakatwhitespace=false,showstringspaces=false,frame=lines,
 escapeinside={(*@}{@*)}}
\title{From an Uninformative Spectral Bound to Useful\\Stopping Guarantees for QC-GHP Codes}
\author{}
\date{October 24, 2026}
\begin{document}
\maketitle
\begin{abstract}
Abiad and Castriota ask for QC-GHP parameters making their spectral stopping-set
bound informative, and propose stronger spectral bounds. A universal
obstruction answers the original useful-parameter question negatively in the
proposition's stated scalar-factor model. Here we show that useful replacements
exist: a degree-corrected fourth-moment criterion gives stopping distance at least
$\gamma+1$, and explicit sparse QC-GHP families meet its hypotheses. A fixed
column/row-weight $4/8$ family has guarantee five; an independently certified
54-variable instance improves on the compared Zhu--Chugg formulas wherever those
formulas apply. A two-block incidence refinement and a separate family with exact
stopping and logical distance six demonstrate further structural gains. Finally,
Gram-equal CSS constructions prove an unbounded limitation for uniformly useful
bounds based only on the specified Gram data. The spectral criterion combines
classical moment and overlap ingredients; the mathematical proofs, constructions and
source-to-matrix bridges are checked in Lean.
\end{abstract}

\section{The source question, its obstruction, and the replacement problem}
\label{sec:question}
Aida Abiad and Nichola Castriota pose the motivating question in Section~6 of
\emph{Spectral and Additive Combinatorial Methods for Cycles and Absorbing Sets
in Lifted-Product Quantum LDPC Codes}, July 15, 2026, arXiv:2607.13666v1 \cite{ref1}.
Can QC-GHP parameters make Proposition~5.1 give an informative stopping-set
bound? Remark~5.2 and Section~6 also motivate tighter spectral bounds.

Throughout, a stopping support of a binary check matrix is a nonempty set of
columns for which no row meets exactly one selected column. Its minimum size is
the stopping distance $s(H)$. This is a property of the specified check matrix.
Logical distance instead minimizes over nontrivial logical words; a lower bound
for logical distance alone need not exclude all small stopping supports.

\subsection{Why the original parameters cannot provide the desired improvement}
The source specializes to a scalar second factor in Section~3.2 and retains that
model for Proposition~5.1. Write the expanded real zero--one matrix as
$H_X=[A\mid I_r\otimes B^{\mathsf T}]$, with $m=r\ell$ rows,
positive column weight $\gamma$, and constant row weight $k$.
Assume positive first-factor column width, $r,\ell>0$, and $m\ge2$ so that
the second singular value is defined. Put $\eta=\sigma_2(H_X)^2$, the second
largest eigenvalue of $H_XH_X^{\mathsf T}$, counting multiplicity. The source
requires the strict condition
\[
 \eta<\gamma/2,
 \qquad a_\pm=\frac m\gamma\left(1\pm\sqrt{1-2\eta/\gamma}\right).
\]
Its large-root conclusion additionally requires excluding the small-size branch:
for a stopping support $S$, one must establish $a_-<|S|$. The spectral inequality
alone is therefore not a useful stopping-distance guarantee.

\begin{priorresult}
For every binary biregular scalar-factor matrix with these positive dimensions
and degrees, it is impossible to have both $\eta<\gamma/2$ and $a_+>2$.
If the strict spectral condition holds, then $r=1$, $1<a_+\le2$, and every
stopping support satisfies $|S|\ge2$ and $a_-<|S|$.
\end{priorresult}

The proof, formalized in \cite{ref7}, separates two unavoidable regimes. For $r\ge2$, the
repeated scalar block forces $\eta\ge\gamma^2$, contradicting the strict
condition. For $r=1$, binary regularity, edge counting and the Gram trace force
the dense regime in which $a_+\le2$ and $a_-<2$. Positive columns already exclude
singleton stopping supports, so the separate small-branch requirement is met
but yields no improvement. The all-ones $2\times4$ example has $\eta=0$ and
stopping distance two: raw spectral-gap examples do exist.
The complete proof is \texttt{SourceQuestionAnswer.lean} and its frozen
\texttt{Spectral.lean} dependency; Appendix~\ref{app:links} gives full current
and immutable URLs for the complete formal proof sources.

\subsection{Useful spectral and incidence replacements}
The obstruction makes the replacement problem substantive: can different
spectral information give a stopping bound greater than two on an explicit
sparse QC-GHP family? Can incidence information give useful guarantees when
cross-block intersections exceed within-block intersections? We answer both
existence questions positively. These are precise replacement questions motivated
by the source direction; optimal universal refinements are a separate target.
\begin{longtable}{@{}>{\raggedright\arraybackslash}p{.465\linewidth}>{\raggedright\arraybackslash}p{.465\linewidth}@{}}
\toprule
\caption{Source question, replacement questions and exact answers.}\label{tab:answers}\\
\textbf{Question / source} & \textbf{Answer and scope}\\\midrule\endhead
Abiad--Castriota, 15 July 2026, v1, Prop. 5.1, Remark 5.2, §6: useful parameters for the stated bound & Negative, by the parameter obstruction (Section 1.1).\\[4pt]
Same source: tighter spectral stopping bounds, including multi-eigenvalue refinements & YES for useful-replacement existence: a fourth-moment criterion proves \ensuremath{\gamma}+1 on explicit sparse joint Sidon families. Theorem 1 and Section 3 prove existence; Theorem 5 limits uniform improvement on other families.\\[4pt]
Incidence-conditioned replacement: a stopping guarantee exceeding two on a sparse QC-GHP family & Positive: Theorems 2--3 give the criterion and families; Theorem 4 gives exact distance six.\\[4pt]
\bottomrule\end{longtable}
The central argument is as follows. Section~\ref{sec:moment} replaces the failed
second-eigenvalue regime by an exact fourth-moment criterion. Section~\ref{sec:families}
realizes it on actual sparse CSS matrices and compares the resulting guarantees
with Zhu--Chugg. Section~\ref{sec:incidence} retains block incidence to sharpen
classical overlap counting; Section~\ref{sec:six} shows how additional finite
incidence information yields an exact distance. Section~\ref{sec:gram} then proves
why the specified Gram data cannot uniformly recover those incidence gains.

\subsection{Reference chronology and calendar-time comparison}
Table~\ref{tab:chronology} makes the central source dates visible before the
proofs. Every elapsed time is measured to this paper's date,
October 24, 2026 (Pacific), rather than to a moving reading date. These are source
ages, not claims about how long a problem was unresolved. Appendix~\ref{app:chronology}
records all used sources and held leads, with version distinctions and precision.
\begin{center}\small
\begin{longtable}{@{}p{.27\linewidth}p{.43\linewidth}p{.22\linewidth}@{}}
\caption{Central sources and elapsed calendar days to October 24, 2026.}\label{tab:chronology}\\
\toprule Source / role & Relevant preprint or inspected milestone & Elapsed days\\\midrule\endhead
Abiad-Castriota & first and inspected v1: 2026-07-15 & 101\\[5pt]
Kashyap--Vardy (overlap bound) & reported conference presentation: 2003\newline reported original composition: 2003-09 & 8333--8697\newline 8425--8454\\[5pt]
Xia-Fu & earliest preprint: 2006-06-12\newline inspected arXivv2 stamp: 2006-11-11 & 7439\newline 7287\\[5pt]
Zhu--Chugg (spectral comparison) & conference predecessor metadata: 2005-09-28\newline inspected journal version: 2007-12 & 7696\newline 6872--6902\\[5pt]
Schwartz-Vardy & first preprint metadata: 2005-03-23\newline inspected journal issue: 2006-03 & 7885\newline 7512--7542\\[5pt]
Wang-Pryadko & first and inspected v1 submission: 2022-03-31\newline PDF internal date: 2022-04-01 & 1668\newline 1667\\[5pt]
Kuo-Ouyang & first preprint metadata: 2024-11-20\newline inspected article-in-press: 2026-03-13 & 703\newline 225\\[5pt]
Dehghan-Banihashemi & first and inspectedv1: 2019-05-31\newline journal volume metadata: 2019 & 2703\newline 2489--2853\\[5pt]
Donovan-Rao-Üsküplü-Yaz\i{}c\i{} & first and inspectedv1: 2022-05-01\newline journal publication year metadata: 2023 & 1637\newline 1028--1392\\[5pt]
Chee-Ling-Yin & first and inspectedv1: 2010-08-10\newline journal issue metadata: 2010-11 & 5919\newline 5807--5836\\[5pt]
Gholami-Samadieh & first and inspectedv1: 2014-05-15\newline internal PDF title-page date: 2022-03-28 & 4545\newline 1671\\[5pt]
Amirzade-Sadeghi & first and inspectedv1: 2017-09-04\newline internal PDF running draft date: 2018-11-18 & 3337\newline 2897\\[5pt]
\bottomrule\end{longtable}\end{center}
For Kashyap--Vardy, the author manuscript is undated; the reported 2003
presentation/composition is distinguished from our inspection. Month- and
year-only milestones give ranges rather than invented publication days.

\section{Replacing the small spectral gap by a fourth-moment criterion}
\label{sec:moment}
The failed regime uses only the second singular value. A different spectral
quantity can certify a concrete incidence property: distinct columns share at
most one check. This property makes a classical stopping-set count useful.

Let $H$ be a finite binary incidence matrix, interpreted over the reals when
forming its Gram matrix. Write $d_v$ for its column degrees, $r_c$ for its row
degrees, and $\eta_1,\ldots,\eta_m$ for all eigenvalues of $HH^{\mathsf T}$,
counting multiplicity. Thus $M_4=\sum_i\eta_i^2$ is its fourth singular-value
moment. All these quantities refer to the actual expanded matrix.

\begin{theorem}[Zero-excess spectral stopping bound]\label{thm:moment}
Suppose every column has degree at least $\gamma$ and
\begin{equation}\label{eq:moment}
 M_4+\sum_v d_v=\sum_v d_v^2+\sum_c r_c^2.
\end{equation}
Every nonempty stopping support has size at least $\gamma+1$.
For a $(\gamma,k)$-regular matrix with $N$ columns and $m$ rows,
condition~\eqref{eq:moment} is equivalently
\[
 M_4+N\gamma=N\gamma^2+mk^2,
 \qquad\text{or}\qquad M_4=N\gamma(\gamma+k-1),
\]
using the edge count $mk=N\gamma$. There is no second-eigenvalue gap
assumption or separate exclusion of small stopping supports.
\end{theorem}
\begin{proof}
Let $q_{uv}$ be the number of checks common to columns $u,v$.
The real spectral identity gives
\[
 M_4=\operatorname{tr}((HH^{\mathsf T})^2)
     =\sum_{u,v}q_{uv}^2,
 \qquad q_{vv}=d_v,\qquad \sum_{u,v}q_{uv}=\sum_cr_c^2.
\]
Consequently the excess in~\eqref{eq:moment} is
\begin{equation}\label{eq:excess}
 \Delta:=M_4+\sum_vd_v-\sum_vd_v^2-\sum_cr_c^2
        =\sum_{u\ne v}q_{uv}(q_{uv}-1).
\end{equation}
Each term is nonnegative because $q_{uv}$ is a natural number. If $\Delta=0$,
then $q_{uv}\le1$ for distinct columns. Choose a column $v$ in a stopping
support $S$. Every one of its at least $\gamma$ checks must meet another
selected column. Each other selected column can account for at most one such
check. Hence $\gamma\le |S|-1$.
\end{proof}

This criterion combines known ingredients. Abiad--Castriota \cite{ref1}, Lemma~4.8,
gives the moment/cycle identity; Kashyap--Vardy \cite{ref2}, Theorem~1, gives the
uniform-overlap count, and Xia--Fu \cite{ref8}, Theorem~3, supplies a related
minimum-column-weight formulation. The contribution here is the checked
spectral corollary, its explicit QC-GHP realizations and the complete
source-to-matrix proof chain. \texttt{SpectralMoment.lean} and
\texttt{MomentBound.lean} check the identities and the stopping implication.

\paragraph{A bound valid outside the informative regime.}
For positive column weight $\gamma$, define $L$ from $(m,N,\gamma,k,M_4)$ to be
$\gamma+1$ when the regular moment equality holds and two otherwise.
\texttt{MomentReplacement.lean} proves $2\le L\le |S|$ for every nonempty
stopping support of every such regular binary matrix. The fallback follows
from positive columns excluding singletons. \texttt{GlobalMomentFamily.lean}
checks that $L=\gamma+1>2$ on the families below for $\gamma\ge2$.
Thus the construction realizes an actual globally valid bound function.

\paragraph{Small positive excess.}
The same calculation has a quantitative extension. For a natural number $\mu$,
if $\Delta<2\mu(\mu+1)$, every distinct-column overlap is at most $\mu$:
an overlap at least $\mu+1$ would contribute at least $2\mu(\mu+1)$ through
the two ordered pairs in~\eqref{eq:excess}. Covering the neighbors of a selected
column then gives
\[
 \gamma+\mu\le\mu|S|.
\]
\texttt{MomentExcess.lean} checks this implication without regularity or an
assumed overlap bound. It is a conditional criterion; we do not claim a new
infinite sparse QC family realizing its useful positive-excess regime.

\section{Explicit sparse families and the useful guarantee}
\label{sec:families}
An abstract criterion alone would leave the existence question unanswered.
We now construct actual QC-GHP matrices satisfying it, including CSS
commutation, degree counts and nontrivial logical content.

\subsection{A uniform construction for every even column weight}
Fix an even $\gamma\ge2$ and an odd $n\ge3^{2\gamma}$. In $\mathbb Z/n\mathbb Z$ set
\[
 S=\{3^i:0\le i<\gamma\},\qquad
 T=\{3^{\gamma+i}:0\le i<\gamma\}.
\]
Let $A,B$ be their binary circulants, with $A_{c,j}=1$ precisely when
$c-j\in S$, and similarly for $B,T$. Use the scalar-factor convention
\begin{equation}\label{eq:joint}
 H_X=[A\mid B^{\mathsf T}],\qquad H_Z=[B\mid A^{\mathsf T}].
\end{equation}
The combined support is Sidon, including repeated summands; all relevant pair
sums are less than $n$. Its two supports are disjoint and distinct columns of
$H_X$ meet in at most one check. Circulants commute, so over $\mathbb F_2$,
$H_XH_Z^{\mathsf T}=AB^{\mathsf T}+B^{\mathsf T}A=0$.
There are $n$ checks and $2n$ variables, column weight $\gamma$ and row weight
$2\gamma$. The actual moment and resulting bound are
\[
 M_4=2n\gamma(3\gamma-1),\qquad s(H_X)\ge\gamma+1.
\]
Theorem~\ref{thm:moment} therefore gives an informative stopping guarantee.
Both sectors also have nonzero kernel weight at least $\gamma+1$.

These matrices have nontrivial logical content. The word equal to one on the
left block and zero on the right has zero syndrome since $\gamma$ is even.
Its left parity is odd since $n$ is odd, while every companion row, and every
linear combination of companion rows, has even left parity. Thus it is outside
the stabilizer span. For fixed $\gamma$, choosing
$n=3^{2\gamma}(2t+1)$ gives arbitrarily long matrices of fixed degrees.
In particular $\gamma=2,n=81$ gives 162 variables and stopping guarantee three.
No encoded-dimension formula or positive-rate claim is made.
\texttt{JointSidon}, \texttt{JointMatrices}, \texttt{JointPowers},
\texttt{JointMomentFamily} and \texttt{JointKernel} check these claims.

\subsection{Short supports: guarantee five at fixed weights $4/8$}
The exponential sufficient length is unnecessary for a useful weight-four
example. For every odd $n\ge37$ use
\begin{equation}\label{eq:short}
 S=\{0,1,3,11\},\qquad T=\{0,5,12,18\}
\end{equation}
in~\eqref{eq:joint}. Then
\[
 N=2n,\quad \gamma=4,\quad k=8,\quad M_4=88n,\quad L=5.
\]
Every nonempty X-stopping support has size at least five; every nonzero kernel
word in either sector has weight at least five. The all-ones-left class is
nontrivial. Thus a fixed-degree infinite sparse family starts at 74 variables.

To verify the incidence hypotheses, it suffices that each support is separately
Sidon and that the mixed sums $s+t$ are unique with their ordered support roles.
Disjoint supports and a Sidon union are unnecessary. All relevant natural pair
sums are at most 36, so reduction modulo $n\ge37$ preserves these properties.
\texttt{SeparateSidon.lean}, \texttt{ShortFamilySupports.lean} and
\texttt{ShortFamily.lean} check the finite identities, no-wrap bridge, actual
matrices, moment, both-sector bounds and logical witness. This realizes the
girth-at-least-six difference criterion of \cite{ref1}, Corollary~4.11; the classical
criterion receives credit rather than being presented as newly discovered.

The same supports also work modulo 27, by a separate finite certificate rather
than the no-wrap argument. \texttt{Practical27Core.lean} checks every
distinct-column overlap of the actual 54-variable matrix in Lean's kernel.
The remaining \texttt{Practical27} modules prove
$M_4=2376$, stopping bound five, both-sector kernel bounds and nontrivial
logical content. A search suggested the supports; no external search outcome
is a proof premise. Neither this instance nor the general short family is
claimed to have exact stopping distance five.

\subsection{Comparison with Zhu--Chugg on the certified matrices}
Zhu--Chugg \cite{ref3}, Corollary~6, equations~41--42, gives the two expressions
\begin{equation}\label{eq:zhu}
 \frac{N(2\gamma-\eta)}{\gamma k-\eta},\qquad
 \frac{2N(2\gamma+k-2-\eta)}{k(\gamma k-\eta)},
\end{equation}
where $\eta$ is the second Gram eigenvalue. Their regular connected setting,
strict gap $\eta<\gamma k$ and other applicability hypotheses must be respected.

For the 54-variable instance, take the real check vectors
\[
 u=\delta_0+\delta_8-\delta_4-\delta_{12},\qquad
 v=\delta_{20}+\delta_1-\delta_{24}-\delta_5.
\]
They have disjoint supports, squared norms four, Gram energies 34 and cross
energy zero. The two-dimensional Rayleigh certificate therefore gives
$\eta\ge17/2$. Under the required gap $\eta<32$, the first expression
in~\eqref{eq:zhu} is negative and the second is at most $297/94<4$.
The new guarantee five is stronger even after integer rounding.
\texttt{Practical27Spectral.lean} checks the actual matrix evaluations,
eigenvalue implication and formula comparison.

For the joint-support example $\gamma=2,n=81$, another exact plane certificate
has squared norms 40, energies 280 and cross energy 38. It gives
$\eta\ge121/20>6$. Under $\eta<8$ both expressions~\eqref{eq:zhu} are negative,
whereas the moment criterion gives three. \texttt{JointSpectralBase.lean}
checks the comparison. These are exact certificates, not floating-point
premises; when a required gap fails, the cited formulas are outside their
stipulated regime.

\section{Retaining block incidence strengthens overlap counting}
\label{sec:incidence}
The moment criterion certifies overlap at most one. A different useful regime
allows larger overlaps between column blocks while retaining smaller overlaps
inside each block. Keeping this distinction sharpens the classical count.

\begin{theorem}[Two-block incidence bound]\label{thm:partition}
Partition the columns of a finite binary incidence relation into two parts.
Let $S$ be a nonempty stopping support. Suppose every selected column has at
least $\gamma$ neighbors, distinct selected columns in the same part share at
most $\sigma$ checks, and selected columns in different parts share at most
$\tau$ checks, where $\sigma\le\tau$. Then
\[
 2(\gamma+\sigma)\le(\sigma+\tau)|S|.
\]
Either part is allowed to be empty within $S$. For $\sigma+\tau>0$, the
resulting integer lower bound is $\lceil2(\gamma+\sigma)/(\sigma+\tau)\rceil$.
\end{theorem}
\begin{proof}
For a selected column $v$, every incident check must meet another selected
column, so
\[
 \gamma\le\sum_{w\in S\setminus\{v\}}|N(v)\cap N(w)|.
\]
Write $p,q$ for the selected counts in the two parts. A selected left column
gives $\gamma+\sigma\le\sigma p+\tau q$; a selected right column gives
$\gamma+\sigma\le\tau p+\sigma q$. When both parts occur, add the inequalities.
If only one occurs, its inequality is $\gamma+\sigma\le\sigma|S|$,
which is at least as strong because $\tau\ge\sigma$.
\end{proof}
\texttt{PartitionBound.lean} checks the cover, finite sums, partition and
one-part cases. The uniform version is $\gamma+\mu\le\mu|S|$.
Kashyap--Vardy \cite{ref2}, Theorem~1, assumes uniform column size $w$ and maximum
intersection $\mu$, giving $1+w/\mu$. At $\gamma=4,\sigma=1,\tau=2$,
its uniform count gives three, while Theorem~\ref{thm:partition} gives four.

\subsection{Sidon supports verify the required intersections}
For a finite abelian group $G$, call $T$ Sidon if equality $a+b=c+d$ among its
elements forces equality of the unordered pairs, including repeated elements.
Let $A_{c,j}=1$ precisely when $c-j\in T$, and set
\[
 H_X=H_Z=[A\mid A^{\mathsf T}].
\]
The two column neighborhoods are $j+T$ and $j-T$. Distinct columns in the same
part share at most one check. Across parts they share at most two: after fixing
one common check $c_0$, the Sidon identity leaves only $c_0$ and $v+w-c_0$.
Every column has degree $|T|$ and every row has degree $2|T|$.
Over $\mathbb F_2$, commutation gives
$H_XH_Z^{\mathsf T}=AA^{\mathsf T}+A^{\mathsf T}A=0$.
These are the scalar QC-GHP/generalized-bicycle matrices of \cite{ref1,ref4}, with
$r_1=n_1=1$ and $b=a$. Thus
\[
 3|S|\ge2(|T|+1).
\]
\texttt{SidonStopping.lean} and \texttt{SidonMatrices.lean} include the actual
matrix, degree and cardinality bridges, with empty-intersection cases.

\subsection{Sparse families for any prescribed fixed guarantee}
\begin{theorem}[Paired powers family]\label{thm:paired}
Fix an even integer $\gamma\ge2$ and an odd $n\ge3^\gamma$. In
$G=\mathbb Z/n\mathbb Z$ take
\[
 T=\{1,3,3^2,\ldots,3^{\gamma-1}\},\qquad
 a(x)=\sum_{i=0}^{\gamma-1}x^{3^i},\qquad H_X=H_Z=[A\mid A^{\mathsf T}].
\]
These are binary CSS checks of column weight $\gamma$ and row weight $2\gamma$,
with stopping distance at least $\lceil2(\gamma+1)/3\rceil$ and a nontrivial
logical class. For fixed $\gamma$, the lengths $n=3^\gamma(2t+1)$ are unbounded
while the degrees stay fixed.
\end{theorem}
\begin{proof}
Sort exponents in $3^a+3^b=3^c+3^d$ so that $a\le b$ and $c\le d$.
If $b<d$, then the left side is at most $2\cdot3^b$, while the right is at
least $3^d\ge3\cdot3^b$, a contradiction. The opposite inequality is symmetric;
cancel equal largest terms to identify the smaller exponents. All relevant pair
sums are below $n$, since $2\cdot3^{\gamma-1}<3^\gamma\le n$. This proves the
Sidon property after reduction and distinctness of the $\gamma$ marks.
Theorem~\ref{thm:partition} applies with $\sigma=1,\tau=2$.
The all-ones-left word is a kernel word of odd left parity, whereas every
companion stabilizer combination has even left parity. It is therefore a
nontrivial logical word.
\end{proof}
\texttt{PowersThree.lean}, \texttt{PowersFamily.lean} and
\texttt{SidonMatrices.lean} check the support, no-wrap, matrix and span claims.
\begin{center}
\begin{tabular}{@{}ccc@{}}\toprule
Column / row weights & Two-block guarantee & Uniform-overlap guarantee\\\midrule
$4/8$ & $4$ & $3$\\
$6/12$ & $5$ & $4$\\
$8/16$ & $6$ & $5$\\\bottomrule
\end{tabular}
\end{center}
Choose $\gamma$ first to obtain any desired fixed guarantee, then increase
length with degrees fixed. This does not establish growing stopping distance at
fixed degree, positive asymptotic rate or practical parameter efficiency.
The sufficient length is exponential in $\gamma$; a companion-row stopping
support gives the upper witness $2\gamma$, checked in \texttt{CSSStopping.lean}.

A second choice, $\gamma=2d$ and $n=3^{2d}$, gives stopping and nonzero-kernel
weight at least $d+1$. \texttt{GrowingFamily.lean} proves
$(2d)3^d\le n$, so incidence density is at most $3^{-d}$;
\texttt{KernelFamily.lean} supplies the kernel bridge. Here the degree varies
with length.

\section{Additional incidence information gives exact distance six}
\label{sec:six}
The reusable criterion need not be sharp. A separate finite certificate and
cover argument can close the gap between it and a logical upper witness.

\begin{theorem}[Exact-six family]\label{thm:six}
For every positive integer $t$, let $n=39t$, let $A$ have circulant support
$\{0,1,4,6\}$, and set $H_X=H_Z=[A\mid A^{\mathsf T}]$.
There are $78t$ variables and $39t$ checks, column weight four and row weight
eight. Both stopping distance and quantum logical distance are exactly six.
\end{theorem}
\begin{proof}
The two-block criterion gives four. To exclude sizes at most five, simultaneous
translation reduces the base $n=39$ check to one root in each side.
A sound degree-one extension procedure rejects both roots with four further
additions. Its exact Boolean tests are reduced in Lean's kernel; no external
search result is a premise.

Reduction modulo 39 maps lifted column lists to their base lists, lifts each
incident base check, and is injective on each check's neighbor set. Thus a
lifted stopping support projects to a base stopping support of no greater size.
The checked local-cover theorem transfers the lower bound six to every
positive multiple of 39.

For the upper bound, select
\[
 \operatorname{left}(0),\ \operatorname{left}(-5),\ \operatorname{left}(-4),\quad
 \operatorname{right}(1),\ \operatorname{right}(5),\ \operatorname{right}(6).
\]
These six columns are distinct at every stated lift. Every incident check has
a selected partner; an independent exact parity calculation shows all selected
check degrees are even, so the word is in the binary kernel. Its left weight
is three, while each companion row has left weight four. Left parity therefore
separates it from the entire stabilizer span. Every nonzero kernel word has
stopping support; the stopping lower bound and this nontrivial six-word give
logical distance six in both coinciding sectors.
\end{proof}
\texttt{BicycleBaseCore}, \texttt{BicycleBase}, \texttt{BicycleFamilyCore} and
\texttt{BicycleFamily} check the finite certificate and cover transfer.
\texttt{BicycleMatrices}, \texttt{BicycleDegrees}, \texttt{BicycleLogical} and
\texttt{KernelStopping} check actual matrices, degrees, syndrome, cardinality
and nonstabilizer status. For the B1 support characterization cited in \cite{ref7}, the minimum witnesses
are companion rows; here the upper witness is a nontrivial logical word. It is not assumed elementary:
the base reconstruction has selected check degrees zero, two and four.
No general formula for encoded dimension is claimed.

For comparison with~\eqref{eq:zhu}, an exact two-dimensional plane certificate
gives $\eta\ge290/13$ on every 39-cover. The base vectors have disjoint supports,
squared norms 13, energies 332 and cross energy 42; the minimum plane quotient
is $(332-42)/13$. The checked fiber argument scales norms and energies exactly
through each cover. Thus both Zhu--Chugg expressions are negative whenever
their required gap $\eta<32$ holds. \texttt{SpectralPlane}, \texttt{FiberSums},
\texttt{BicycleCoverEnergy} and \texttt{BicycleSpectralFamily} check this
comparison without assuming a floating-point spectral envelope.

\section{Why the full check Gram data cannot uniformly recover the gain}
\label{sec:gram}
The positive results suggest a limitation question: could a universally valid
function of the full Gram data recover these incidence gains on every family?
The following indistinguishable good/bad pairs obstruct that possibility for
the specified inputs.

\begin{theorem}[Restricted Gram-data obstruction]\label{thm:gram}
For a constituent circulant $A$ from the paired families above, form
\[
 \begin{aligned}
 \text{Good: }& H_X=H_Z=[A\mid A^{\mathsf T}],\\
 \text{Bad: }& H_X'=[A\mid A],\qquad H_Z'=[A^{\mathsf T}\mid A^{\mathsf T}].
 \end{aligned}
\]
Both are binary CSS pairs with the same check-matrix dimensions and degrees
and the same two separate real check Gram matrices. The bad code has stopping
distance two and a nontrivial logical word of weight two. Any real-valued
stopping lower-bound function depending only on these shared data and valid
on both members is at most two on the good member as well.
\end{theorem}
\begin{proof}
Circulants commute with their transposes. Over the reals,
\[
 H_XH_X^{\mathsf T}=AA^{\mathsf T}+A^{\mathsf T}A
                   =2AA^{\mathsf T}=H_X'(H_X')^{\mathsf T},
\]
and the same calculation gives equality for the Z check Grams. Both pairs
commute in the binary CSS sense. The bad X-matrix has identical positive
columns at each coordinate, giving a stopping pair; positive columns exclude
singletons. The corresponding weight-two word is outside the companion
stabilizer span, as checked by the binary parity separation in
\texttt{BadLogical.lean}. The analogous sector statement is also checked.
Validity of a bound on the bad pair forces its value to be at most two;
the identical inputs force the same output on the good pair.
\end{proof}

For an unbounded separation, choose $\gamma=4M$ for $M\ge1$ and
$n=3^{4M}(2t+1)$ in Theorem~\ref{thm:paired}. Every good stopping support
has size at least $2M+1$, while its bad twin has stopping distance two.
Consequently any such bound $L$ satisfies $ML<|S|$ on every good stopping
support. \texttt{RealGramBound.lean}, \texttt{GramGap.lean} and
\texttt{BothGram.lean} check the exact quantifiers and both Gram identities;
the chain includes \texttt{GramObstruction}, \texttt{BadStopping},
\texttt{SidonBad} and \texttt{BadDegrees}.

This obstruction includes functions of the complete Gram spectra and of
eigenvectors determined solely by those check Gram matrices. It does not
assume equality of binary ranks or encoded dimensions, and it does not include
cross-sector real incidence or variable-side eigenvectors as inputs.
It is compatible with Theorem~\ref{thm:moment}: the zero-excess criterion succeeds
on a particular class, whereas the paired-family obstruction exhibits a
different class on which shared Gram data cannot distinguish large from small
stopping distance.

\section{Relation to classical stopping sets and quantum-code results}
\label{sec:related}
The classical ingredients explain both the positive answer and its scope.
Kashyap--Vardy \cite{ref2} already gives a uniform-overlap route to a guarantee above
two in the weight-four regime. Theorem~\ref{thm:partition} improves its value
by retaining the partition. Xia--Fu \cite{ref8}, Theorem~3, assumes minimum column
weight $\gamma$, maximum distinct-column intersection $\lambda$ and integral
$\gamma/\lambda$, and gives pseudo-weight at least $\gamma/\lambda+1$ with an
equality characterization. At $\gamma=4,\lambda=2$ this is three, compared
with four from the partition criterion. Their girth-at-least-six Theorem~1
does not apply to the paired family's cross-block four-cycles.

Schwartz--Vardy \cite{ref5} emphasizes that stopping distance belongs to a specified
parity-check matrix and that codeword supports are stopping sets. Their
stopping redundancy optimizes over added checks, unlike the fixed-$H_X$
guarantees here. Wang--Pryadko \cite{ref4}, equation~8, uses
$H_X=[A\mid B]$, $H_Z^{\mathsf T}=[B;-A]$; our paired case takes
$B=A^{\mathsf T}$. Their logical-distance bounds and dimension formula concern
different invariants. Likewise, Kuo--Ouyang \cite{ref6} studies decoding exploiting
degeneracy; decoder success does not establish an ordinary peeling
stopping-distance bound.

Dehghan--Banihashemi \cite{ref9} gives cospectral bipartite examples with matched
degree sequences and different cycle counts. Those cycle statements explain
how spectral data can omit incidence information; they do not supply the
actual full-Gram-equal CSS pairs and stopping/logical certificates proved here.
The existing source screen is reused without counting a new paper cohort.

Donovan--Rao--\"Usk\"upl\"u--Yaz\i c\i\ \cite{ref10}, Lemma~9, gives stopping
distance at least eight when $3\nmid a$ for its specified DM/DCA matrices and
four-row-group hypotheses. No QC-GHP CSS-companion bridge for our displayed
matrices is supplied by that statement. Chee--Ling--Yin \cite{ref11} distinguishes
cyclic difference packing from the partition condition for PCDP; our
overlapping supports are not claimed to be a PCDP. Gholami--Samadieh \cite{ref12}
studies classical lifted finite set systems and inevitable walks, and
Amirzade--Sadeghi \cite{ref13} gives classical single/multiple-edge difference criteria
and lifting bounds. These comparisons do not make classical stopping bounds
above two new. Search coverage and access limitations are recorded once in
Appendix~\ref{app:ledger} and the accompanying source ledgers.

\section{Answers, implications, and remaining targets}
\label{sec:answers}
The original parameter question is answered negatively by the parameter
obstruction: Proposition~5.1's strict spectral regime cannot produce a large-root
guarantee above two in its stated model, with the small branch properly handled.
The stronger-replacement existence question is answered positively by
Theorem~\ref{thm:moment} and the explicit sparse families. The
incidence-conditioned existence question is answered positively by
Theorems~\ref{thm:partition}--\ref{thm:paired}; Theorem~\ref{thm:six}
additionally supplies exact stopping and logical distance six on a separate
infinite family. These are the answers, complete mathematical statements with their displayed hypotheses.

The implications have distinct scopes. Full-spectrum moment information is
sufficient on a concrete class. Retaining block incidence improves uniform
overlap counting. Even the full separate check Gram matrices cannot uniformly
recover the incidence gains on the class in Theorem~\ref{thm:gram}.
Binary ranks, encoded dimensions and cross-sector data remain outside that
obstruction. The short $4/8$ family has lower bound five, not an exact-five
theorem, and no positive asymptotic rate or growing distance at fixed degree
has been established.

Further mathematical questions are to construct or obstruct informative sparse scalar
QC families with positive moment excess, or determine sharp stopping distances
for the general powers families. The two replacement existence questions are
settled by the displayed constructions.

\paragraph{Formal evidence.}
The mathematical proofs and source-to-matrix bridges are checked in Lean.
The build and audit cover 485 named declarations in the replacement-bound
proof project; the parameter-obstruction sources have their own audits.
Appendix~\ref{app:proofs} maps the claims to their complete proof chains,
Appendix~\ref{app:links} gives every full current and immutable Lean URL, and
Appendix~\ref{app:chronology} gives the complete chronology behind
Table~\ref{tab:chronology}. Appendix~\ref{app:ledger} states the bounded
prior-art coverage; detailed search and verification records accompany the sources.

\begin{thebibliography}{99}
\bibitem{ref1} Aida Abiad and Nichola Castriota. Spectral and Additive Combinatorial Methods for Cycles and Absorbing Sets in Lifted-Product Quantum LDPC Codes. 15 July 2026, v1, Prop. 5.1, Remark 5.2,§6. \url{https://arxiv.org/html/2607.13666v1}.

\bibitem{ref2} Navin Kashyap and Alexander Vardy. Stopping Sets in Codes from Designs. Undated revised author manuscript; reports composition September 2003 and ISIT 2003 presentation. Theorem 1. \url{https://ece.iisc.ac.in/~nkashyap/Papers/stopsets.pdf}.

\bibitem{ref3} Mingrui Zhu and Keith M.Chugg. Bounds on the Expansion Properties of Tanner Graphs. IEEE TIT 53(12), 4834--4840, December 2007. Corollary 6. \url{https://hal.usc.edu/chugg/docs/pubs/zhu2007bounds.pdf}.

\bibitem{ref4} Renyu Wang and Leonid P.Pryadko. Distance bounds for generalized bicycle codes. arXiv:2203.17216v1,31 March 2022; internal PDF date1 April 2022. \url{https://arxiv.org/abs/2203.17216v1}.

\bibitem{ref5} Moshe Schwartz and Alexander Vardy. On the Stopping Distance and the Stopping Redundancy of Codes. IEEE TIT 52(3), 922--932, March 2006. \url{https://www.ece.mcmaster.ca/~schwam10/publications/MSjrnl6.pdf}.

\bibitem{ref6} Kao-Yueh Kuo and Yingkai Ouyang. Degenerate quantum erasure decoding. npj Quantum Information 12, 75, 2026. Inspected article-in-press main PDF; version distinctions in chronology. \url{https://doi.org/10.1038/s41534-026-01212-3}.

\bibitem{ref7} Formal proof sources for the scalar-factor parameter obstruction and the B1 stopping-support characterization. \url{https://github.com/alectionary/vibe-math/tree/master/formalization/stopping-bound}.

\bibitem{ref8} Shu-Tao Xia and Fang-Wei Fu. Minimum Pseudo-Weight and Minimum Pseudo-Codewords of LDPC Codes. arXiv:cs/0606051v2,11 November 2006; first submission12 June 2006. Internal PDF date2019 distinguished in chronology. \url{https://arxiv.org/abs/cs/0606051v2}.

\bibitem{ref9} Ali Dehghan and Amir H.Banihashemi. Cospectral Bipartite Graphs with the Same Degree Sequences but with Different Number of Large Cycles. arXiv:1905.13228v1,31 May 2019; existing full-text screen reused. \url{https://arxiv.org/abs/1905.13228v1}.

\bibitem{ref10} Diane Donovan,Asha Rao,Elif Üsküplü,E.\c{S}ule Yaz\i{}c\i{}. QC-LDPC Codes from Difference Matrices and Difference Covering Arrays. arXiv:2205.00563v1, 1 May 2022. Lemma 9,Section 3,Section 5. \url{https://arxiv.org/abs/2205.00563v1}.

\bibitem{ref11} Yeow Meng Chee,Alan C.H.Ling,Jianxing Yin. Optimal Partitioned Cyclic Difference Packings for Frequency Hopping and Code Synchronization. arXiv:1008.1615v1, 10 August 2010. Section II,Theorem 6.3. \url{https://arxiv.org/abs/1008.1615v1}.

\bibitem{ref12} Mohammad Gholami,Mehdi Samadieh. Quasi Cyclic LDPC Codes Based on Finite Set Systems. arXiv:1405.3775v1, 15 May 2014; internal PDF date 28 March 2022 separately recorded. \url{https://arxiv.org/abs/1405.3775v1}.

\bibitem{ref13} Farzane Amirzade,Mohammad-Reza Sadeghi. Lower bounds on the lifting degree of single-edge and multiple-edge QC-LDPC codes by difference matrices. arXiv:1709.00825v1, 4 September 2017; internal running draft date 18 November 2018. Full-document target screen, no additional explicit unresolved target. \url{https://arxiv.org/abs/1709.00825v1}.

\end{thebibliography}
\appendix
\clearpage\section{Proof coverage and checked Lean excerpts}\label{app:proofs}
The aggregate build and all 485 named-declaration audits pass in Lean 4.36.0-rc1 with mathlib \nolinkurl{5f7b4198d22a9c82e1c698a379334fe128095c82}. Only propext, Classical.choice and Quot.sound occur. The project contains 50 proof modules and four imported research modules. Reproduce with lake build StrongerAnswer, followed by lake env lean AuditStronger.lean. The logs, signatures, complete source links, hashes and independent reconstructions accompany the archive.

\begin{longtable}{@{}>{\raggedright\arraybackslash}p{.465\linewidth}>{\raggedright\arraybackslash}p{.465\linewidth}@{}}
\toprule
\textbf{Proof component} & \textbf{Checked source}\\\midrule\endhead
Ordered eigenvalue fourth moment, zero excess and useful spectral bound & SpectralMoment; MomentBound\\[4pt]
Globally valid bound function and quantitative nonzero excess & MomentReplacement; GlobalMomentFamily; MomentExcess\\[4pt]
Shorter fixed-degree family and 54-variable matrix, both sectors & SeparateSidon; ShortFamilySupports; ShortFamily; Practical27Core; Practical27; Practical27Design; Practical27Spectral\\[4pt]
Actual source matrices, length/density consumers and bad companion degrees & SourceConsumers; BadDegrees\\[4pt]
Constructed joint Sidon matrices, spectral family and both binary sectors & JointSidon; JointMatrices; JointPowers; JointMomentFamily; JointKernel\\[4pt]
Exact plane certificates and all 39-cover energy scaling & SpectralPlane; FiberSums; BicycleSpectralBase; BicycleCoverEnergy; BicycleSpectralFamily; JointSpectralBase\\[4pt]
Real-valued/unbounded Gram gap and nontrivial bad logical class & RealGramBound; GramGap; BadLogical; BothGram\\[4pt]
Growing-degree density and kernel guarantees & GrowingFamily; KernelFamily\\[4pt]
Generic neighbor cover and two-block/uniform bounds & PartitionBound.lean\\[4pt]
Sidon intersection limits and no-wrap powers-of-three support & SidonStopping.lean; PowersThree.lean\\[4pt]
Binary matrices, CSS, degree counts and logical separation & SidonMatrices.lean; PowersFamily.lean\\[4pt]
Rooted search soundness, two roots and cyclic covers & dependencies/IncidenceCore,List,Cover; BicycleBaseCore,Base,FamilyCore,Family\\[4pt]
Concrete actual matrices, degrees, overlap and exact logical witness & BicycleMatrices,Degrees,Overlap,Logical; KernelStopping\\[4pt]
Same-Gram identity, bad pair and universal-bound obstruction & GramObstruction; BadStopping; SidonBad\\[4pt]
Known nonzero companion-row stopping mechanism & CSSStopping.lean\\[4pt]
Consumers and exhaustive named-declaration audit & StrongerAnswer.lean; AuditStronger.lean\\[4pt]
\bottomrule\end{longtable}
\subsection*{\texttt{PartitionBound.lean}}
\begin{lstlisting}
theorem partition_bound (color : V (*@\ensuremath{\to}@*) Bool) (S : Finset V)
    (hS : Incidence.Stopping E S) ((*@\ensuremath{\gamma}@*) (*@\ensuremath{\sigma}@*) (*@\ensuremath{\tau}@*) : (*@\ensuremath{\mathbb N}@*)) (hst : (*@\ensuremath{\sigma}@*) (*@\ensuremath{\le}@*) (*@\ensuremath{\tau}@*))
    (hdeg : (*@\ensuremath{\forall}@*) v (*@\ensuremath{\in}@*) S, (*@\ensuremath{\gamma}@*) (*@\ensuremath{\le}@*) (neighbors E v).card)
    (hov : (*@\ensuremath{\forall}@*) v (*@\ensuremath{\in}@*) S, (*@\ensuremath{\forall}@*) w (*@\ensuremath{\in}@*) S, v (*@\ensuremath{\ne}@*) w (*@\ensuremath{\to}@*)
      overlap E v w (*@\ensuremath{\le}@*) if color v = color w then (*@\ensuremath{\sigma}@*) else (*@\ensuremath{\tau}@*)) :
    2 * ((*@\ensuremath{\gamma}@*) + (*@\ensuremath{\sigma}@*)) (*@\ensuremath{\le}@*) ((*@\ensuremath{\sigma}@*) + (*@\ensuremath{\tau}@*)) * S.card := by
  classical
  let P := S.filter (fun w => color w = false)
  let Q := S.filter (fun w => color w = true)
  have hcard : P.card + Q.card = S.card := by
    have h := card_filter_add_card_filter_not (s := S) (fun w => color w = false)
    simpa [P,Q,Bool.not_eq_false] using h
  have hfalse (v : V) (hv : v (*@\ensuremath{\in}@*) S) (hc : color v = false) :
      (*@\ensuremath{\gamma}@*) + (*@\ensuremath{\sigma}@*) (*@\ensuremath{\le}@*) (*@\ensuremath{\sigma}@*) * P.card + (*@\ensuremath{\tau}@*) * Q.card := by
    have h := column_budget E color S hS (*@\ensuremath{\gamma}@*) (*@\ensuremath{\sigma}@*) (*@\ensuremath{\tau}@*) hdeg hov v hv
    simpa [hc,P,Q,Bool.not_eq_false] using h
  have htrue (v : V) (hv : v (*@\ensuremath{\in}@*) S) (hc : color v = true) :
      (*@\ensuremath{\gamma}@*) + (*@\ensuremath{\sigma}@*) (*@\ensuremath{\le}@*) (*@\ensuremath{\tau}@*) * P.card + (*@\ensuremath{\sigma}@*) * Q.card := by
    have h := column_budget E color S hS (*@\ensuremath{\gamma}@*) (*@\ensuremath{\sigma}@*) (*@\ensuremath{\tau}@*) hdeg hov v hv
    simpa [hc,P,Q,Bool.not_eq_true,add_comm] using h
  by_cases hp : P.Nonempty
  (*@\ensuremath{\cdot}@*) obtain (*@\ensuremath{\langle}@*)v,hv(*@\ensuremath{\rangle}@*) := hp
    have hv' := mem_filter.mp hv
    have hb := hfalse v hv'.1 hv'.2
    by_cases hq : Q.Nonempty
    (*@\ensuremath{\cdot}@*) obtain (*@\ensuremath{\langle}@*)w,hw(*@\ensuremath{\rangle}@*) := hq
      have hw' := mem_filter.mp hw
      have hb' := htrue w hw'.1 hw'.2
      nlinarith
    (*@\ensuremath{\cdot}@*) have hq0 : Q.card = 0 := card_eq_zero.mpr (not_nonempty_iff_eq_empty.mp hq)
      nlinarith [Nat.mul_le_mul_right P.card hst]
  (*@\ensuremath{\cdot}@*) have hp0 : P.card = 0 := card_eq_zero.mpr (not_nonempty_iff_eq_empty.mp hp)
    obtain (*@\ensuremath{\langle}@*)v,hv(*@\ensuremath{\rangle}@*) := hS.1
    have hc : color v = true := by
      cases hc : color v
      (*@\ensuremath{\cdot}@*) have : v (*@\ensuremath{\in}@*) P := mem_filter.mpr (*@\ensuremath{\langle}@*)hv,hc(*@\ensuremath{\rangle}@*)
        exact False.elim (hp (*@\ensuremath{\langle}@*)v,this(*@\ensuremath{\rangle}@*))
      (*@\ensuremath{\cdot}@*) rfl
    have hb := htrue v hv hc
    nlinarith [Nat.mul_le_mul_right Q.card hst]
\end{lstlisting}
\subsection*{\texttt{PowersFamily.lean}}
\begin{lstlisting}
theorem positive_family ((*@\ensuremath{\gamma}@*) n : (*@\ensuremath{\mathbb N}@*)) [NeZero n]
    (hn : 3^(*@\ensuremath{\gamma}@*) (*@\ensuremath{\le}@*) n) (heven : Even (*@\ensuremath{\gamma}@*)) (hodd : Odd n) :
    let S := support (*@\ensuremath{\gamma}@*) n
    (SidonStopping.X (ZMod 2) S * (SidonStopping.X (ZMod 2) S)(*@\ensuremath{{}^{\mathsf T}}@*) = 0) (*@\ensuremath{\land}@*)
    ((*@\ensuremath{\forall}@*) v, (Finset.univ.filter (fun c => SidonStopping.X (ZMod 2) S c v = 1)).card = (*@\ensuremath{\gamma}@*)) (*@\ensuremath{\land}@*)
    ((*@\ensuremath{\forall}@*) c, (Finset.univ.filter (fun v => SidonStopping.X (ZMod 2) S c v = 1)).card = 2*(*@\ensuremath{\gamma}@*)) (*@\ensuremath{\land}@*)
    ((*@\ensuremath{\forall}@*) T, Incidence.Stopping (SidonStopping.E S) T (*@\ensuremath{\to}@*) 2*((*@\ensuremath{\gamma}@*)+1) (*@\ensuremath{\le}@*) 3*T.card) (*@\ensuremath{\land}@*)
    (SidonStopping.X (ZMod 2) S).mulVec (SidonStopping.leftOnes (G := ZMod n)) = 0 (*@\ensuremath{\land}@*)
    ((*@\ensuremath{\neg}@*) (*@\ensuremath{\exists}@*) a, (SidonStopping.X (ZMod 2) S)(*@\ensuremath{{}^{\mathsf T}}@*).mulVec a =
      SidonStopping.leftOnes (G := ZMod n)) (*@\ensuremath{\land}@*)
    (SidonStopping.X (*@\ensuremath{\mathbb R}@*) S * (SidonStopping.X (*@\ensuremath{\mathbb R}@*) S)(*@\ensuremath{{}^{\mathsf T}}@*) =
      SidonStopping.BadX (*@\ensuremath{\mathbb R}@*) S * (SidonStopping.BadX (*@\ensuremath{\mathbb R}@*) S)(*@\ensuremath{{}^{\mathsf T}}@*)) := by
  dsimp only
  have hcard := support_card (*@\ensuremath{\gamma}@*) n hn
  have he : Even (support (*@\ensuremath{\gamma}@*) n).card := by simpa [hcard] using heven
  refine (*@\ensuremath{\langle}@*)SidonStopping.css _,?_,?_,?_,SidonStopping.leftOnes_kernel _ he,?_,SidonStopping.gram_equal _(*@\ensuremath{\rangle}@*)
  (*@\ensuremath{\cdot}@*) intro v; rw [SidonStopping.column_weight,hcard]
  (*@\ensuremath{\cdot}@*) intro c; rw [SidonStopping.row_weight,hcard]
  (*@\ensuremath{\cdot}@*) intro T hT
    simpa only [hcard] using SidonStopping.stopping_bound _ (sidon (*@\ensuremath{\gamma}@*) n hn) T hT
  (*@\ensuremath{\cdot}@*) apply SidonStopping.leftOnes_not_stabilizer _ he
    simpa using hodd
\end{lstlisting}
\subsection*{\texttt{KernelStopping.lean}}
\begin{lstlisting}
theorem quantum_distance_six (h : 39 (*@\ensuremath{\mid}@*) n) :
    (HX (n := n) (ZMod 2) * (HZ (n := n) (ZMod 2))(*@\ensuremath{{}^{\mathsf T}}@*) = 0) (*@\ensuremath{\land}@*)
    ((*@\ensuremath{\forall}@*) x : Var n (*@\ensuremath{\to}@*) ZMod 2, x (*@\ensuremath{\ne}@*) 0 (*@\ensuremath{\to}@*) (HX (ZMod 2)).mulVec x = 0 (*@\ensuremath{\to}@*)
      6 (*@\ensuremath{\le}@*) (BinaryStopping.support x).card) (*@\ensuremath{\land}@*)
    (HX (ZMod 2)).mulVec (@word n) = 0 (*@\ensuremath{\land}@*)
    ((*@\ensuremath{\neg}@*) (*@\ensuremath{\exists}@*) a : Check n (*@\ensuremath{\to}@*) ZMod 2, ((HZ (ZMod 2))(*@\ensuremath{{}^{\mathsf T}}@*).mulVec a) = (@word n)) (*@\ensuremath{\land}@*)
    (BinaryStopping.support (@word n)).card = 6 :=
  (*@\ensuremath{\langle}@*)css,all_nonzero_kernel_weight_ge_six h,word_kernel h,word_not_stabilizer h,word_weight h(*@\ensuremath{\rangle}@*)
\end{lstlisting}
\subsection*{\texttt{SidonBad.lean}}
\begin{lstlisting}
theorem gram_only_bound (S : Finset G)
    (lower : Matrix G G (*@\ensuremath{\mathbb R}@*) (*@\ensuremath{\to}@*) (*@\ensuremath{\mathbb N}@*))
    (valid_bad : (*@\ensuremath{\forall}@*) T : Finset (Bool (*@\ensuremath{\times}@*) G), Incidence.Stopping (badE S) T (*@\ensuremath{\to}@*)
      lower (BadX (*@\ensuremath{\mathbb R}@*) S * (BadX (*@\ensuremath{\mathbb R}@*) S)(*@\ensuremath{{}^{\mathsf T}}@*)) (*@\ensuremath{\le}@*) T.card) :
    lower (X (*@\ensuremath{\mathbb R}@*) S * (X (*@\ensuremath{\mathbb R}@*) S)(*@\ensuremath{{}^{\mathsf T}}@*)) (*@\ensuremath{\le}@*) 2 := by
  rw [gram_equal]
  have h := valid_bad twin (twin_stopping S)
  simpa [twin_card] using h
\end{lstlisting}
\subsection*{\texttt{GramGap.lean}}
\begin{lstlisting}
theorem unbounded_gram_gap (M t : (*@\ensuremath{\mathbb N}@*)) [NeZero (liftLength (2*M) t)]
    (lower : Matrix (ZMod (liftLength (2*M) t)) (ZMod (liftLength (2*M) t)) (*@\ensuremath{\mathbb R}@*) (*@\ensuremath{\to}@*) (*@\ensuremath{\mathbb R}@*))
    (valid_bad : (*@\ensuremath{\forall}@*) U : Finset (Bool (*@\ensuremath{\times}@*) ZMod (liftLength (2*M) t)),
      Incidence.Stopping (SidonStopping.badE (support (4*M) (liftLength (2*M) t))) U (*@\ensuremath{\to}@*)
      lower (SidonStopping.BadX (*@\ensuremath{\mathbb R}@*) (support (4*M) (liftLength (2*M) t)) *
        (SidonStopping.BadX (*@\ensuremath{\mathbb R}@*) (support (4*M) (liftLength (2*M) t)))(*@\ensuremath{{}^{\mathsf T}}@*)) (*@\ensuremath{\le}@*) (U.card : (*@\ensuremath{\mathbb R}@*)))
    (T : Finset (Bool (*@\ensuremath{\times}@*) ZMod (liftLength (2*M) t)))
    (hT : Incidence.Stopping (SidonStopping.E (support (4*M) (liftLength (2*M) t))) T) :
    (M:(*@\ensuremath{\mathbb R}@*))*lower (SidonStopping.X (*@\ensuremath{\mathbb R}@*) (support (4*M) (liftLength (2*M) t)) *
      (SidonStopping.X (*@\ensuremath{\mathbb R}@*) (support (4*M) (liftLength (2*M) t)))(*@\ensuremath{{}^{\mathsf T}}@*)) < (T.card : (*@\ensuremath{\mathbb R}@*)) := by
  have hbound := SidonStopping.real_gram_only_bound _ lower valid_bad
  have hrewrite : 4*M = 2*(2*M) := by omega
  rw [hrewrite] at hT
  have hg := arbitrary_fixed_guarantee (2*M) t T hT
  have hgr : (2*M+1:(*@\ensuremath{\mathbb N}@*)) (*@\ensuremath{\le}@*) T.card := hg
  have hrr : (2*(M:(*@\ensuremath{\mathbb R}@*))+1) (*@\ensuremath{\le}@*) (T.card:(*@\ensuremath{\mathbb R}@*)) := by exact_mod_cast hgr
  have hmul := mul_le_mul_of_nonneg_left hbound (Nat.cast_nonneg M : (0:(*@\ensuremath{\mathbb R}@*)) (*@\ensuremath{\le}@*) M)
  linarith
\end{lstlisting}
\subsection*{\texttt{BicycleSpectralFamily.lean}}
\begin{lstlisting}
theorem zhu_chugg_negative (h : 39 (*@\ensuremath{\mid}@*) n)
    (hgap : eigs (Family.HX (n := n) (*@\ensuremath{\mathbb R}@*)) (*@\ensuremath{\langle}@*)1,by exact check_count h(*@\ensuremath{\rangle}@*) < 32) :
    let s := eigs (Family.HX (n := n) (*@\ensuremath{\mathbb R}@*)) (*@\ensuremath{\langle}@*)1,by exact check_count h(*@\ensuremath{\rangle}@*)
    (2*(n:(*@\ensuremath{\mathbb R}@*)))*(8-s)/(32-s) < 0 (*@\ensuremath{\land}@*)
    (4*(n:(*@\ensuremath{\mathbb R}@*)))*(14-s)/(8*(32-s)) < 0 := by
  dsimp
  have hs := second_lower h
  have hn : (0:(*@\ensuremath{\mathbb R}@*)) < n := by exact_mod_cast NeZero.pos n
  constructor
  (*@\ensuremath{\cdot}@*) apply div_neg_of_neg_of_pos
    (*@\ensuremath{\cdot}@*) exact mul_neg_of_pos_of_neg (by positivity) (by linarith)
    (*@\ensuremath{\cdot}@*) linarith
  (*@\ensuremath{\cdot}@*) apply div_neg_of_neg_of_pos
    (*@\ensuremath{\cdot}@*) exact mul_neg_of_pos_of_neg (by positivity) (by linarith)
    (*@\ensuremath{\cdot}@*) positivity
\end{lstlisting}
\subsection*{\texttt{MomentBound.lean}}
\begin{lstlisting}
theorem regular_spectral_bound (E : V (*@\ensuremath{\to}@*) C (*@\ensuremath{\to}@*) Prop) [DecidableRel E] ((*@\ensuremath{\gamma}@*) k : (*@\ensuremath{\mathbb N}@*))
    (hcol : (*@\ensuremath{\forall}@*) v, (PartitionStopping.neighbors E v).card = (*@\ensuremath{\gamma}@*))
    (hrow : (*@\ensuremath{\forall}@*) c, rowWeight E c = k)
    (hmoment : ((*@\ensuremath{\sum}@*) i,(StoppingSpectral.eigs (H E) i)^2) + (Fintype.card V : (*@\ensuremath{\mathbb R}@*))*(*@\ensuremath{\gamma}@*) =
      (Fintype.card V : (*@\ensuremath{\mathbb R}@*))*((*@\ensuremath{\gamma}@*):(*@\ensuremath{\mathbb R}@*))^2 + (Fintype.card C : (*@\ensuremath{\mathbb R}@*))*(k:(*@\ensuremath{\mathbb R}@*))^2)
    (S : Finset V) (hS : Incidence.Stopping E S) : (*@\ensuremath{\gamma}@*)+1 (*@\ensuremath{\le}@*) S.card := by
  apply spectral_stopping_bound E (*@\ensuremath{\gamma}@*) (fun v => by rw [hcol]) _ S hS
  simpa only [hcol,hrow,sum_const,card_univ,nsmul_eq_mul] using hmoment
\end{lstlisting}
\subsection*{\texttt{JointMomentFamily.lean}}
\begin{lstlisting}
theorem fourth_moment (hn : 3^(2*(*@\ensuremath{\gamma}@*)) (*@\ensuremath{\le}@*) n) :
    ((*@\ensuremath{\sum}@*) i,(eigs (JointSidon.X (*@\ensuremath{\mathbb R}@*) (leftMarks (*@\ensuremath{\gamma}@*) n) (rightMarks (*@\ensuremath{\gamma}@*) n)) i)^2) =
      2*(n:(*@\ensuremath{\mathbb R}@*))*((*@\ensuremath{\gamma}@*):(*@\ensuremath{\mathbb R}@*))*(3*((*@\ensuremath{\gamma}@*):(*@\ensuremath{\mathbb R}@*))-1) := by
  have hm := StoppingMoment.moment_of_overlap_one
    (JointSidon.E (leftMarks (*@\ensuremath{\gamma}@*) n) (rightMarks (*@\ensuremath{\gamma}@*) n))
    (JointSidon.overlap_one _ _ (joint_sidon (*@\ensuremath{\gamma}@*) n hn) (disjoint_marks (*@\ensuremath{\gamma}@*) n hn))
  rw [real_matrix] at hm
  simp_rw [column_degree hn,row_degree hn] at hm
  simp only [sum_const,card_univ,Fintype.card_prod,Fintype.card_bool,ZMod.card,
    Nat.cast_mul,Nat.cast_ofNat,nsmul_eq_mul] at hm
  nlinarith
\end{lstlisting}
\subsection*{\texttt{ShortFamily.lean}}
\begin{lstlisting}
theorem source_stopping_five (hn : 37(*@\ensuremath{\le}@*)n) (U : Finset (Bool (*@\ensuremath{\times}@*) ZMod n))
    (hU : Incidence.Stopping (fun v c => HX (n := n) (ZMod 2) c v=1) U) : 5(*@\ensuremath{\le}@*)U.card := by
  refine regular_spectral_bound (E (n := n)) 4 8 (degree hn) (row_degree hn) ?_ U ?_
  (*@\ensuremath{\cdot}@*) rw [real_matrix,fourth_moment hn]
    simp only [Fintype.card_prod,Fintype.card_bool,ZMod.card,Nat.cast_mul,Nat.cast_ofNat]
    ring
  (*@\ensuremath{\cdot}@*) simpa only [HX,Incidence.Stopping,JointSidon.binary_incidence] using hU

theorem replacement_value (hn : 37(*@\ensuremath{\le}@*)n) :
    regularReplacement n (2*n) 4 8 ((*@\ensuremath{\sum}@*) i,(eigs (HX (n := n) (*@\ensuremath{\mathbb R}@*)) i)^2)=5 := by
  rw [fourth_moment hn]
  unfold regularReplacement
  split_ifs with he
  (*@\ensuremath{\cdot}@*) norm_num
  (*@\ensuremath{\cdot}@*) exfalso
    apply he
    push_cast
    ring
\end{lstlisting}
\subsection*{\texttt{Practical27Spectral.lean}}
\begin{lstlisting}
theorem zhu_chugg_comparison (hgap : eigs HR (*@\ensuremath{\langle}@*)1,by decide(*@\ensuremath{\rangle}@*) < 32) :
    let s := eigs HR (*@\ensuremath{\langle}@*)1,by decide(*@\ensuremath{\rangle}@*)
    54*(8-s)/(32-s) < 0 (*@\ensuremath{\land}@*)
      108*(14-s)/(8*(32-s)) (*@\ensuremath{\le}@*) (297:(*@\ensuremath{\mathbb R}@*))/94 (*@\ensuremath{\land}@*) (297:(*@\ensuremath{\mathbb R}@*))/94 < 4 := by
  dsimp
  have hs := second_lower
  refine (*@\ensuremath{\langle}@*)?_,?_,by norm_num(*@\ensuremath{\rangle}@*)
  (*@\ensuremath{\cdot}@*) apply div_neg_of_neg_of_pos
    (*@\ensuremath{\cdot}@*) nlinarith
    (*@\ensuremath{\cdot}@*) linarith
  (*@\ensuremath{\cdot}@*) apply (div_le_iff(*@\ensuremath{{}_0}@*) (by positivity : 0 < 8*(32-eigs HR (*@\ensuremath{\langle}@*)1,by decide(*@\ensuremath{\rangle}@*)))).mpr
    nlinarith
\end{lstlisting}
\paragraph{Dependency guide.} PartitionBound provides generic counting. SidonStopping and PowersThree prove incidence hypotheses constructively. SidonMatrices and PowersFamily supply actual binary CSS and nontrivial logical content. BicycleBaseCore/List dependencies give search soundness; BicycleBase rejects small roots; FamilyCore/Cover transfer exact6. BicycleMatrices/Degrees/Logical/KernelStopping check the actual code and logical witness. GramObstruction, BadStopping and SidonBad prove indistinguishability and exact bad stopping2. CSSStopping checks the known stabilizer upper mechanism. AuditStronger checks every named declaration, including every frozen dependency. SpectralMoment/MomentBound establish the ordered-eigenvalue criterion. JointSidon/Matrices/Powers/MomentFamily/Kernel construct the useful spectral families. JointSpectralBase checks the cited-formula comparison. RealGramBound/GramGap/BadLogical establish the real-valued Gram obstruction and nontrivial bad logical words. SpectralPlane/FiberSums and the cover-energy modules check exact two-vector certificates and cover transfer. GrowingFamily/KernelFamily give the stopping/kernel growth guarantees. SeparateSidon/ShortFamilySupports/ShortFamily provide shorter fixed-degree parameters; Practical27Core/Practical27/Practical27Design/Practical27Spectral check the 54-variable example and exact cited-formula comparison. MomentReplacement/GlobalMomentFamily and MomentExcess provide the global bound function and quantitative extension. BothGram covers both separate check Gram matrices, excluding cross-sector/binary-rank information.
\clearpage\section{Complete current and immutable Lean URLs}\label{app:links}
Every URL below is printed in full and clickable. The replacement-bound sources are pinned at commit \texttt{bed7719cb3cac4118e4f333690db364566ede03e}. The parameter-obstruction and B1 proof sources have separate immutable pins.

\subsection{This paper: every proof, audit and imported research file}
\medskip
\par\noindent\textbf{\texttt{AuditStronger.lean}}\par\nopagebreak
{\footnotesize\noindent Current: \url{https://github.com/alectionary/vibe-math/blob/master/formalization/stopping-bound/stronger/AuditStronger.lean}\\
Immutable: \url{https://github.com/alectionary/vibe-math/blob/bed7719cb3cac4118e4f333690db364566ede03e/formalization/stopping-bound/stronger/AuditStronger.lean}\par}
\smallskip
\medskip
\par\noindent\textbf{\texttt{BadDegrees.lean}}\par\nopagebreak
{\footnotesize\noindent Current: \url{https://github.com/alectionary/vibe-math/blob/master/formalization/stopping-bound/stronger/BadDegrees.lean}\\
Immutable: \url{https://github.com/alectionary/vibe-math/blob/bed7719cb3cac4118e4f333690db364566ede03e/formalization/stopping-bound/stronger/BadDegrees.lean}\par}
\smallskip
\medskip
\par\noindent\textbf{\texttt{BadLogical.lean}}\par\nopagebreak
{\footnotesize\noindent Current: \url{https://github.com/alectionary/vibe-math/blob/master/formalization/stopping-bound/stronger/BadLogical.lean}\\
Immutable: \url{https://github.com/alectionary/vibe-math/blob/bed7719cb3cac4118e4f333690db364566ede03e/formalization/stopping-bound/stronger/BadLogical.lean}\par}
\smallskip
\medskip
\par\noindent\textbf{\texttt{BadStopping.lean}}\par\nopagebreak
{\footnotesize\noindent Current: \url{https://github.com/alectionary/vibe-math/blob/master/formalization/stopping-bound/stronger/BadStopping.lean}\\
Immutable: \url{https://github.com/alectionary/vibe-math/blob/bed7719cb3cac4118e4f333690db364566ede03e/formalization/stopping-bound/stronger/BadStopping.lean}\par}
\smallskip
\medskip
\par\noindent\textbf{\texttt{BicycleBase.lean}}\par\nopagebreak
{\footnotesize\noindent Current: \url{https://github.com/alectionary/vibe-math/blob/master/formalization/stopping-bound/stronger/BicycleBase.lean}\\
Immutable: \url{https://github.com/alectionary/vibe-math/blob/bed7719cb3cac4118e4f333690db364566ede03e/formalization/stopping-bound/stronger/BicycleBase.lean}\par}
\smallskip
\medskip
\par\noindent\textbf{\texttt{BicycleBaseCore.lean}}\par\nopagebreak
{\footnotesize\noindent Current: \url{https://github.com/alectionary/vibe-math/blob/master/formalization/stopping-bound/stronger/BicycleBaseCore.lean}\\
Immutable: \url{https://github.com/alectionary/vibe-math/blob/bed7719cb3cac4118e4f333690db364566ede03e/formalization/stopping-bound/stronger/BicycleBaseCore.lean}\par}
\smallskip
\medskip
\par\noindent\textbf{\texttt{BicycleCoverEnergy.lean}}\par\nopagebreak
{\footnotesize\noindent Current: \url{https://github.com/alectionary/vibe-math/blob/master/formalization/stopping-bound/stronger/BicycleCoverEnergy.lean}\\
Immutable: \url{https://github.com/alectionary/vibe-math/blob/bed7719cb3cac4118e4f333690db364566ede03e/formalization/stopping-bound/stronger/BicycleCoverEnergy.lean}\par}
\smallskip
\medskip
\par\noindent\textbf{\texttt{BicycleDegrees.lean}}\par\nopagebreak
{\footnotesize\noindent Current: \url{https://github.com/alectionary/vibe-math/blob/master/formalization/stopping-bound/stronger/BicycleDegrees.lean}\\
Immutable: \url{https://github.com/alectionary/vibe-math/blob/bed7719cb3cac4118e4f333690db364566ede03e/formalization/stopping-bound/stronger/BicycleDegrees.lean}\par}
\smallskip
\medskip
\par\noindent\textbf{\texttt{BicycleFamily.lean}}\par\nopagebreak
{\footnotesize\noindent Current: \url{https://github.com/alectionary/vibe-math/blob/master/formalization/stopping-bound/stronger/BicycleFamily.lean}\\
Immutable: \url{https://github.com/alectionary/vibe-math/blob/bed7719cb3cac4118e4f333690db364566ede03e/formalization/stopping-bound/stronger/BicycleFamily.lean}\par}
\smallskip
\medskip
\par\noindent\textbf{\texttt{BicycleFamilyCore.lean}}\par\nopagebreak
{\footnotesize\noindent Current: \url{https://github.com/alectionary/vibe-math/blob/master/formalization/stopping-bound/stronger/BicycleFamilyCore.lean}\\
Immutable: \url{https://github.com/alectionary/vibe-math/blob/bed7719cb3cac4118e4f333690db364566ede03e/formalization/stopping-bound/stronger/BicycleFamilyCore.lean}\par}
\smallskip
\medskip
\par\noindent\textbf{\texttt{BicycleLogical.lean}}\par\nopagebreak
{\footnotesize\noindent Current: \url{https://github.com/alectionary/vibe-math/blob/master/formalization/stopping-bound/stronger/BicycleLogical.lean}\\
Immutable: \url{https://github.com/alectionary/vibe-math/blob/bed7719cb3cac4118e4f333690db364566ede03e/formalization/stopping-bound/stronger/BicycleLogical.lean}\par}
\smallskip
\medskip
\par\noindent\textbf{\texttt{BicycleMatrices.lean}}\par\nopagebreak
{\footnotesize\noindent Current: \url{https://github.com/alectionary/vibe-math/blob/master/formalization/stopping-bound/stronger/BicycleMatrices.lean}\\
Immutable: \url{https://github.com/alectionary/vibe-math/blob/bed7719cb3cac4118e4f333690db364566ede03e/formalization/stopping-bound/stronger/BicycleMatrices.lean}\par}
\smallskip
\medskip
\par\noindent\textbf{\texttt{BicycleOverlap.lean}}\par\nopagebreak
{\footnotesize\noindent Current: \url{https://github.com/alectionary/vibe-math/blob/master/formalization/stopping-bound/stronger/BicycleOverlap.lean}\\
Immutable: \url{https://github.com/alectionary/vibe-math/blob/bed7719cb3cac4118e4f333690db364566ede03e/formalization/stopping-bound/stronger/BicycleOverlap.lean}\par}
\smallskip
\medskip
\par\noindent\textbf{\texttt{BicycleSpectralBase.lean}}\par\nopagebreak
{\footnotesize\noindent Current: \url{https://github.com/alectionary/vibe-math/blob/master/formalization/stopping-bound/stronger/BicycleSpectralBase.lean}\\
Immutable: \url{https://github.com/alectionary/vibe-math/blob/bed7719cb3cac4118e4f333690db364566ede03e/formalization/stopping-bound/stronger/BicycleSpectralBase.lean}\par}
\smallskip
\medskip
\par\noindent\textbf{\texttt{BicycleSpectralFamily.lean}}\par\nopagebreak
{\footnotesize\noindent Current: \url{https://github.com/alectionary/vibe-math/blob/master/formalization/stopping-bound/stronger/BicycleSpectralFamily.lean}\\
Immutable: \url{https://github.com/alectionary/vibe-math/blob/bed7719cb3cac4118e4f333690db364566ede03e/formalization/stopping-bound/stronger/BicycleSpectralFamily.lean}\par}
\smallskip
\medskip
\par\noindent\textbf{\texttt{BothGram.lean}}\par\nopagebreak
{\footnotesize\noindent Current: \url{https://github.com/alectionary/vibe-math/blob/master/formalization/stopping-bound/stronger/BothGram.lean}\\
Immutable: \url{https://github.com/alectionary/vibe-math/blob/bed7719cb3cac4118e4f333690db364566ede03e/formalization/stopping-bound/stronger/BothGram.lean}\par}
\smallskip
\medskip
\par\noindent\textbf{\texttt{CSSStopping.lean}}\par\nopagebreak
{\footnotesize\noindent Current: \url{https://github.com/alectionary/vibe-math/blob/master/formalization/stopping-bound/stronger/CSSStopping.lean}\\
Immutable: \url{https://github.com/alectionary/vibe-math/blob/bed7719cb3cac4118e4f333690db364566ede03e/formalization/stopping-bound/stronger/CSSStopping.lean}\par}
\smallskip
\medskip
\par\noindent\textbf{\texttt{FiberSums.lean}}\par\nopagebreak
{\footnotesize\noindent Current: \url{https://github.com/alectionary/vibe-math/blob/master/formalization/stopping-bound/stronger/FiberSums.lean}\\
Immutable: \url{https://github.com/alectionary/vibe-math/blob/bed7719cb3cac4118e4f333690db364566ede03e/formalization/stopping-bound/stronger/FiberSums.lean}\par}
\smallskip
\medskip
\par\noindent\textbf{\texttt{GlobalMomentFamily.lean}}\par\nopagebreak
{\footnotesize\noindent Current: \url{https://github.com/alectionary/vibe-math/blob/master/formalization/stopping-bound/stronger/GlobalMomentFamily.lean}\\
Immutable: \url{https://github.com/alectionary/vibe-math/blob/bed7719cb3cac4118e4f333690db364566ede03e/formalization/stopping-bound/stronger/GlobalMomentFamily.lean}\par}
\smallskip
\medskip
\par\noindent\textbf{\texttt{GramGap.lean}}\par\nopagebreak
{\footnotesize\noindent Current: \url{https://github.com/alectionary/vibe-math/blob/master/formalization/stopping-bound/stronger/GramGap.lean}\\
Immutable: \url{https://github.com/alectionary/vibe-math/blob/bed7719cb3cac4118e4f333690db364566ede03e/formalization/stopping-bound/stronger/GramGap.lean}\par}
\smallskip
\medskip
\par\noindent\textbf{\texttt{GramObstruction.lean}}\par\nopagebreak
{\footnotesize\noindent Current: \url{https://github.com/alectionary/vibe-math/blob/master/formalization/stopping-bound/stronger/GramObstruction.lean}\\
Immutable: \url{https://github.com/alectionary/vibe-math/blob/bed7719cb3cac4118e4f333690db364566ede03e/formalization/stopping-bound/stronger/GramObstruction.lean}\par}
\smallskip
\medskip
\par\noindent\textbf{\texttt{GrowingFamily.lean}}\par\nopagebreak
{\footnotesize\noindent Current: \url{https://github.com/alectionary/vibe-math/blob/master/formalization/stopping-bound/stronger/GrowingFamily.lean}\\
Immutable: \url{https://github.com/alectionary/vibe-math/blob/bed7719cb3cac4118e4f333690db364566ede03e/formalization/stopping-bound/stronger/GrowingFamily.lean}\par}
\smallskip
\medskip
\par\noindent\textbf{\texttt{JointKernel.lean}}\par\nopagebreak
{\footnotesize\noindent Current: \url{https://github.com/alectionary/vibe-math/blob/master/formalization/stopping-bound/stronger/JointKernel.lean}\\
Immutable: \url{https://github.com/alectionary/vibe-math/blob/bed7719cb3cac4118e4f333690db364566ede03e/formalization/stopping-bound/stronger/JointKernel.lean}\par}
\smallskip
\medskip
\par\noindent\textbf{\texttt{JointMatrices.lean}}\par\nopagebreak
{\footnotesize\noindent Current: \url{https://github.com/alectionary/vibe-math/blob/master/formalization/stopping-bound/stronger/JointMatrices.lean}\\
Immutable: \url{https://github.com/alectionary/vibe-math/blob/bed7719cb3cac4118e4f333690db364566ede03e/formalization/stopping-bound/stronger/JointMatrices.lean}\par}
\smallskip
\medskip
\par\noindent\textbf{\texttt{JointMomentFamily.lean}}\par\nopagebreak
{\footnotesize\noindent Current: \url{https://github.com/alectionary/vibe-math/blob/master/formalization/stopping-bound/stronger/JointMomentFamily.lean}\\
Immutable: \url{https://github.com/alectionary/vibe-math/blob/bed7719cb3cac4118e4f333690db364566ede03e/formalization/stopping-bound/stronger/JointMomentFamily.lean}\par}
\smallskip
\medskip
\par\noindent\textbf{\texttt{JointPowers.lean}}\par\nopagebreak
{\footnotesize\noindent Current: \url{https://github.com/alectionary/vibe-math/blob/master/formalization/stopping-bound/stronger/JointPowers.lean}\\
Immutable: \url{https://github.com/alectionary/vibe-math/blob/bed7719cb3cac4118e4f333690db364566ede03e/formalization/stopping-bound/stronger/JointPowers.lean}\par}
\smallskip
\medskip
\par\noindent\textbf{\texttt{JointSidon.lean}}\par\nopagebreak
{\footnotesize\noindent Current: \url{https://github.com/alectionary/vibe-math/blob/master/formalization/stopping-bound/stronger/JointSidon.lean}\\
Immutable: \url{https://github.com/alectionary/vibe-math/blob/bed7719cb3cac4118e4f333690db364566ede03e/formalization/stopping-bound/stronger/JointSidon.lean}\par}
\smallskip
\medskip
\par\noindent\textbf{\texttt{JointSpectralBase.lean}}\par\nopagebreak
{\footnotesize\noindent Current: \url{https://github.com/alectionary/vibe-math/blob/master/formalization/stopping-bound/stronger/JointSpectralBase.lean}\\
Immutable: \url{https://github.com/alectionary/vibe-math/blob/bed7719cb3cac4118e4f333690db364566ede03e/formalization/stopping-bound/stronger/JointSpectralBase.lean}\par}
\smallskip
\medskip
\par\noindent\textbf{\texttt{KernelFamily.lean}}\par\nopagebreak
{\footnotesize\noindent Current: \url{https://github.com/alectionary/vibe-math/blob/master/formalization/stopping-bound/stronger/KernelFamily.lean}\\
Immutable: \url{https://github.com/alectionary/vibe-math/blob/bed7719cb3cac4118e4f333690db364566ede03e/formalization/stopping-bound/stronger/KernelFamily.lean}\par}
\smallskip
\medskip
\par\noindent\textbf{\texttt{KernelStopping.lean}}\par\nopagebreak
{\footnotesize\noindent Current: \url{https://github.com/alectionary/vibe-math/blob/master/formalization/stopping-bound/stronger/KernelStopping.lean}\\
Immutable: \url{https://github.com/alectionary/vibe-math/blob/bed7719cb3cac4118e4f333690db364566ede03e/formalization/stopping-bound/stronger/KernelStopping.lean}\par}
\smallskip
\medskip
\par\noindent\textbf{\texttt{MomentBound.lean}}\par\nopagebreak
{\footnotesize\noindent Current: \url{https://github.com/alectionary/vibe-math/blob/master/formalization/stopping-bound/stronger/MomentBound.lean}\\
Immutable: \url{https://github.com/alectionary/vibe-math/blob/bed7719cb3cac4118e4f333690db364566ede03e/formalization/stopping-bound/stronger/MomentBound.lean}\par}
\smallskip
\medskip
\par\noindent\textbf{\texttt{MomentExcess.lean}}\par\nopagebreak
{\footnotesize\noindent Current: \url{https://github.com/alectionary/vibe-math/blob/master/formalization/stopping-bound/stronger/MomentExcess.lean}\\
Immutable: \url{https://github.com/alectionary/vibe-math/blob/bed7719cb3cac4118e4f333690db364566ede03e/formalization/stopping-bound/stronger/MomentExcess.lean}\par}
\smallskip
\medskip
\par\noindent\textbf{\texttt{MomentReplacement.lean}}\par\nopagebreak
{\footnotesize\noindent Current: \url{https://github.com/alectionary/vibe-math/blob/master/formalization/stopping-bound/stronger/MomentReplacement.lean}\\
Immutable: \url{https://github.com/alectionary/vibe-math/blob/bed7719cb3cac4118e4f333690db364566ede03e/formalization/stopping-bound/stronger/MomentReplacement.lean}\par}
\smallskip
\medskip
\par\noindent\textbf{\texttt{PartitionBound.lean}}\par\nopagebreak
{\footnotesize\noindent Current: \url{https://github.com/alectionary/vibe-math/blob/master/formalization/stopping-bound/stronger/PartitionBound.lean}\\
Immutable: \url{https://github.com/alectionary/vibe-math/blob/bed7719cb3cac4118e4f333690db364566ede03e/formalization/stopping-bound/stronger/PartitionBound.lean}\par}
\smallskip
\medskip
\par\noindent\textbf{\texttt{PowersFamily.lean}}\par\nopagebreak
{\footnotesize\noindent Current: \url{https://github.com/alectionary/vibe-math/blob/master/formalization/stopping-bound/stronger/PowersFamily.lean}\\
Immutable: \url{https://github.com/alectionary/vibe-math/blob/bed7719cb3cac4118e4f333690db364566ede03e/formalization/stopping-bound/stronger/PowersFamily.lean}\par}
\smallskip
\medskip
\par\noindent\textbf{\texttt{PowersThree.lean}}\par\nopagebreak
{\footnotesize\noindent Current: \url{https://github.com/alectionary/vibe-math/blob/master/formalization/stopping-bound/stronger/PowersThree.lean}\\
Immutable: \url{https://github.com/alectionary/vibe-math/blob/bed7719cb3cac4118e4f333690db364566ede03e/formalization/stopping-bound/stronger/PowersThree.lean}\par}
\smallskip
\medskip
\par\noindent\textbf{\texttt{Practical27.lean}}\par\nopagebreak
{\footnotesize\noindent Current: \url{https://github.com/alectionary/vibe-math/blob/master/formalization/stopping-bound/stronger/Practical27.lean}\\
Immutable: \url{https://github.com/alectionary/vibe-math/blob/bed7719cb3cac4118e4f333690db364566ede03e/formalization/stopping-bound/stronger/Practical27.lean}\par}
\smallskip
\medskip
\par\noindent\textbf{\texttt{Practical27Core.lean}}\par\nopagebreak
{\footnotesize\noindent Current: \url{https://github.com/alectionary/vibe-math/blob/master/formalization/stopping-bound/stronger/Practical27Core.lean}\\
Immutable: \url{https://github.com/alectionary/vibe-math/blob/bed7719cb3cac4118e4f333690db364566ede03e/formalization/stopping-bound/stronger/Practical27Core.lean}\par}
\smallskip
\medskip
\par\noindent\textbf{\texttt{Practical27Design.lean}}\par\nopagebreak
{\footnotesize\noindent Current: \url{https://github.com/alectionary/vibe-math/blob/master/formalization/stopping-bound/stronger/Practical27Design.lean}\\
Immutable: \url{https://github.com/alectionary/vibe-math/blob/bed7719cb3cac4118e4f333690db364566ede03e/formalization/stopping-bound/stronger/Practical27Design.lean}\par}
\smallskip
\medskip
\par\noindent\textbf{\texttt{Practical27Spectral.lean}}\par\nopagebreak
{\footnotesize\noindent Current: \url{https://github.com/alectionary/vibe-math/blob/master/formalization/stopping-bound/stronger/Practical27Spectral.lean}\\
Immutable: \url{https://github.com/alectionary/vibe-math/blob/bed7719cb3cac4118e4f333690db364566ede03e/formalization/stopping-bound/stronger/Practical27Spectral.lean}\par}
\smallskip
\medskip
\par\noindent\textbf{\texttt{RealGramBound.lean}}\par\nopagebreak
{\footnotesize\noindent Current: \url{https://github.com/alectionary/vibe-math/blob/master/formalization/stopping-bound/stronger/RealGramBound.lean}\\
Immutable: \url{https://github.com/alectionary/vibe-math/blob/bed7719cb3cac4118e4f333690db364566ede03e/formalization/stopping-bound/stronger/RealGramBound.lean}\par}
\smallskip
\medskip
\par\noindent\textbf{\texttt{SeparateSidon.lean}}\par\nopagebreak
{\footnotesize\noindent Current: \url{https://github.com/alectionary/vibe-math/blob/master/formalization/stopping-bound/stronger/SeparateSidon.lean}\\
Immutable: \url{https://github.com/alectionary/vibe-math/blob/bed7719cb3cac4118e4f333690db364566ede03e/formalization/stopping-bound/stronger/SeparateSidon.lean}\par}
\smallskip
\medskip
\par\noindent\textbf{\texttt{ShortFamily.lean}}\par\nopagebreak
{\footnotesize\noindent Current: \url{https://github.com/alectionary/vibe-math/blob/master/formalization/stopping-bound/stronger/ShortFamily.lean}\\
Immutable: \url{https://github.com/alectionary/vibe-math/blob/bed7719cb3cac4118e4f333690db364566ede03e/formalization/stopping-bound/stronger/ShortFamily.lean}\par}
\smallskip
\medskip
\par\noindent\textbf{\texttt{ShortFamilySupports.lean}}\par\nopagebreak
{\footnotesize\noindent Current: \url{https://github.com/alectionary/vibe-math/blob/master/formalization/stopping-bound/stronger/ShortFamilySupports.lean}\\
Immutable: \url{https://github.com/alectionary/vibe-math/blob/bed7719cb3cac4118e4f333690db364566ede03e/formalization/stopping-bound/stronger/ShortFamilySupports.lean}\par}
\smallskip
\medskip
\par\noindent\textbf{\texttt{SidonBad.lean}}\par\nopagebreak
{\footnotesize\noindent Current: \url{https://github.com/alectionary/vibe-math/blob/master/formalization/stopping-bound/stronger/SidonBad.lean}\\
Immutable: \url{https://github.com/alectionary/vibe-math/blob/bed7719cb3cac4118e4f333690db364566ede03e/formalization/stopping-bound/stronger/SidonBad.lean}\par}
\smallskip
\medskip
\par\noindent\textbf{\texttt{SidonMatrices.lean}}\par\nopagebreak
{\footnotesize\noindent Current: \url{https://github.com/alectionary/vibe-math/blob/master/formalization/stopping-bound/stronger/SidonMatrices.lean}\\
Immutable: \url{https://github.com/alectionary/vibe-math/blob/bed7719cb3cac4118e4f333690db364566ede03e/formalization/stopping-bound/stronger/SidonMatrices.lean}\par}
\smallskip
\medskip
\par\noindent\textbf{\texttt{SidonStopping.lean}}\par\nopagebreak
{\footnotesize\noindent Current: \url{https://github.com/alectionary/vibe-math/blob/master/formalization/stopping-bound/stronger/SidonStopping.lean}\\
Immutable: \url{https://github.com/alectionary/vibe-math/blob/bed7719cb3cac4118e4f333690db364566ede03e/formalization/stopping-bound/stronger/SidonStopping.lean}\par}
\smallskip
\medskip
\par\noindent\textbf{\texttt{SourceConsumers.lean}}\par\nopagebreak
{\footnotesize\noindent Current: \url{https://github.com/alectionary/vibe-math/blob/master/formalization/stopping-bound/stronger/SourceConsumers.lean}\\
Immutable: \url{https://github.com/alectionary/vibe-math/blob/bed7719cb3cac4118e4f333690db364566ede03e/formalization/stopping-bound/stronger/SourceConsumers.lean}\par}
\smallskip
\medskip
\par\noindent\textbf{\texttt{SpectralMoment.lean}}\par\nopagebreak
{\footnotesize\noindent Current: \url{https://github.com/alectionary/vibe-math/blob/master/formalization/stopping-bound/stronger/SpectralMoment.lean}\\
Immutable: \url{https://github.com/alectionary/vibe-math/blob/bed7719cb3cac4118e4f333690db364566ede03e/formalization/stopping-bound/stronger/SpectralMoment.lean}\par}
\smallskip
\medskip
\par\noindent\textbf{\texttt{SpectralPlane.lean}}\par\nopagebreak
{\footnotesize\noindent Current: \url{https://github.com/alectionary/vibe-math/blob/master/formalization/stopping-bound/stronger/SpectralPlane.lean}\\
Immutable: \url{https://github.com/alectionary/vibe-math/blob/bed7719cb3cac4118e4f333690db364566ede03e/formalization/stopping-bound/stronger/SpectralPlane.lean}\par}
\smallskip
\medskip
\par\noindent\textbf{\texttt{StrongerAnswer.lean}}\par\nopagebreak
{\footnotesize\noindent Current: \url{https://github.com/alectionary/vibe-math/blob/master/formalization/stopping-bound/stronger/StrongerAnswer.lean}\\
Immutable: \url{https://github.com/alectionary/vibe-math/blob/bed7719cb3cac4118e4f333690db364566ede03e/formalization/stopping-bound/stronger/StrongerAnswer.lean}\par}
\smallskip
\medskip
\par\noindent\textbf{\texttt{dependencies/IncidenceCore.lean}}\par\nopagebreak
{\footnotesize\noindent Current: \url{https://github.com/alectionary/vibe-math/blob/master/formalization/stopping-bound/stronger/dependencies/IncidenceCore.lean}\\
Immutable: \url{https://github.com/alectionary/vibe-math/blob/bed7719cb3cac4118e4f333690db364566ede03e/formalization/stopping-bound/stronger/dependencies/IncidenceCore.lean}\par}
\smallskip
\medskip
\par\noindent\textbf{\texttt{dependencies/IncidenceCover.lean}}\par\nopagebreak
{\footnotesize\noindent Current: \url{https://github.com/alectionary/vibe-math/blob/master/formalization/stopping-bound/stronger/dependencies/IncidenceCover.lean}\\
Immutable: \url{https://github.com/alectionary/vibe-math/blob/bed7719cb3cac4118e4f333690db364566ede03e/formalization/stopping-bound/stronger/dependencies/IncidenceCover.lean}\par}
\smallskip
\medskip
\par\noindent\textbf{\texttt{dependencies/IncidenceList.lean}}\par\nopagebreak
{\footnotesize\noindent Current: \url{https://github.com/alectionary/vibe-math/blob/master/formalization/stopping-bound/stronger/dependencies/IncidenceList.lean}\\
Immutable: \url{https://github.com/alectionary/vibe-math/blob/bed7719cb3cac4118e4f333690db364566ede03e/formalization/stopping-bound/stronger/dependencies/IncidenceList.lean}\par}
\smallskip
\medskip
\par\noindent\textbf{\texttt{dependencies/Spectral.lean}}\par\nopagebreak
{\footnotesize\noindent Current: \url{https://github.com/alectionary/vibe-math/blob/master/formalization/stopping-bound/stronger/dependencies/Spectral.lean}\\
Immutable: \url{https://github.com/alectionary/vibe-math/blob/bed7719cb3cac4118e4f333690db364566ede03e/formalization/stopping-bound/stronger/dependencies/Spectral.lean}\par}
\smallskip
\subsection{Parameter-obstruction and B1 proof sources}
\medskip
\par\noindent\textbf{\texttt{Audit.lean}}\par\nopagebreak
{\footnotesize\noindent Current: \url{https://github.com/alectionary/vibe-math/blob/master/formalization/stopping-bound/Audit.lean}\\
Immutable: \url{https://github.com/alectionary/vibe-math/blob/b7628643bb256bf9c1a319f16f6e19608116cff6/formalization/stopping-bound/Audit.lean}\par}
\smallskip
\medskip
\par\noindent\textbf{\texttt{Spectral.lean}}\par\nopagebreak
{\footnotesize\noindent Current: \url{https://github.com/alectionary/vibe-math/blob/master/formalization/stopping-bound/Spectral.lean}\\
Immutable: \url{https://github.com/alectionary/vibe-math/blob/b7628643bb256bf9c1a319f16f6e19608116cff6/formalization/stopping-bound/Spectral.lean}\par}
\smallskip
\medskip
\par\noindent\textbf{\texttt{classification/AuditClassification.lean}}\par\nopagebreak
{\footnotesize\noindent Current: \url{https://github.com/alectionary/vibe-math/blob/master/formalization/stopping-bound/classification/AuditClassification.lean}\\
Immutable: \url{https://github.com/alectionary/vibe-math/blob/b7628643bb256bf9c1a319f16f6e19608116cff6/formalization/stopping-bound/classification/AuditClassification.lean}\par}
\smallskip
\medskip
\par\noindent\textbf{\texttt{classification/Classification.lean}}\par\nopagebreak
{\footnotesize\noindent Current: \url{https://github.com/alectionary/vibe-math/blob/master/formalization/stopping-bound/classification/Classification.lean}\\
Immutable: \url{https://github.com/alectionary/vibe-math/blob/b7628643bb256bf9c1a319f16f6e19608116cff6/formalization/stopping-bound/classification/Classification.lean}\par}
\smallskip
\medskip
\par\noindent\textbf{\texttt{classification/ClassificationBranches.lean}}\par\nopagebreak
{\footnotesize\noindent Current: \url{https://github.com/alectionary/vibe-math/blob/master/formalization/stopping-bound/classification/ClassificationBranches.lean}\\
Immutable: \url{https://github.com/alectionary/vibe-math/blob/b7628643bb256bf9c1a319f16f6e19608116cff6/formalization/stopping-bound/classification/ClassificationBranches.lean}\par}
\smallskip
\medskip
\par\noindent\textbf{\texttt{classification/ClassificationCertificate.lean}}\par\nopagebreak
{\footnotesize\noindent Current: \url{https://github.com/alectionary/vibe-math/blob/master/formalization/stopping-bound/classification/ClassificationCertificate.lean}\\
Immutable: \url{https://github.com/alectionary/vibe-math/blob/b7628643bb256bf9c1a319f16f6e19608116cff6/formalization/stopping-bound/classification/ClassificationCertificate.lean}\par}
\smallskip
\medskip
\par\noindent\textbf{\texttt{classification/ClassificationData.lean}}\par\nopagebreak
{\footnotesize\noindent Current: \url{https://github.com/alectionary/vibe-math/blob/master/formalization/stopping-bound/classification/ClassificationData.lean}\\
Immutable: \url{https://github.com/alectionary/vibe-math/blob/b7628643bb256bf9c1a319f16f6e19608116cff6/formalization/stopping-bound/classification/ClassificationData.lean}\par}
\smallskip
\medskip
\par\noindent\textbf{\texttt{classification/ClassificationKernel.lean}}\par\nopagebreak
{\footnotesize\noindent Current: \url{https://github.com/alectionary/vibe-math/blob/master/formalization/stopping-bound/classification/ClassificationKernel.lean}\\
Immutable: \url{https://github.com/alectionary/vibe-math/blob/b7628643bb256bf9c1a319f16f6e19608116cff6/formalization/stopping-bound/classification/ClassificationKernel.lean}\par}
\smallskip
\medskip
\par\noindent\textbf{\texttt{classification/ClassificationSearch.lean}}\par\nopagebreak
{\footnotesize\noindent Current: \url{https://github.com/alectionary/vibe-math/blob/master/formalization/stopping-bound/classification/ClassificationSearch.lean}\\
Immutable: \url{https://github.com/alectionary/vibe-math/blob/b7628643bb256bf9c1a319f16f6e19608116cff6/formalization/stopping-bound/classification/ClassificationSearch.lean}\par}
\smallskip
\medskip
\par\noindent\textbf{\texttt{classification/dependencies/Incidence.lean}}\par\nopagebreak
{\footnotesize\noindent Current: \url{https://github.com/alectionary/vibe-math/blob/master/formalization/stopping-bound/classification/dependencies/Incidence.lean}\\
Immutable: \url{https://github.com/alectionary/vibe-math/blob/b7628643bb256bf9c1a319f16f6e19608116cff6/formalization/stopping-bound/classification/dependencies/Incidence.lean}\par}
\smallskip
\medskip
\par\noindent\textbf{\texttt{classification/dependencies/IncidenceCSS.lean}}\par\nopagebreak
{\footnotesize\noindent Current: \url{https://github.com/alectionary/vibe-math/blob/master/formalization/stopping-bound/classification/dependencies/IncidenceCSS.lean}\\
Immutable: \url{https://github.com/alectionary/vibe-math/blob/b7628643bb256bf9c1a319f16f6e19608116cff6/formalization/stopping-bound/classification/dependencies/IncidenceCSS.lean}\par}
\smallskip
\medskip
\par\noindent\textbf{\texttt{classification/dependencies/IncidenceCertificate.lean}}\par\nopagebreak
{\footnotesize\noindent Current: \url{https://github.com/alectionary/vibe-math/blob/master/formalization/stopping-bound/classification/dependencies/IncidenceCertificate.lean}\\
Immutable: \url{https://github.com/alectionary/vibe-math/blob/b7628643bb256bf9c1a319f16f6e19608116cff6/formalization/stopping-bound/classification/dependencies/IncidenceCertificate.lean}\par}
\smallskip
\medskip
\par\noindent\textbf{\texttt{classification/dependencies/IncidenceCore.lean}}\par\nopagebreak
{\footnotesize\noindent Current: \url{https://github.com/alectionary/vibe-math/blob/master/formalization/stopping-bound/classification/dependencies/IncidenceCore.lean}\\
Immutable: \url{https://github.com/alectionary/vibe-math/blob/b7628643bb256bf9c1a319f16f6e19608116cff6/formalization/stopping-bound/classification/dependencies/IncidenceCore.lean}\par}
\smallskip
\medskip
\par\noindent\textbf{\texttt{classification/dependencies/IncidenceCover.lean}}\par\nopagebreak
{\footnotesize\noindent Current: \url{https://github.com/alectionary/vibe-math/blob/master/formalization/stopping-bound/classification/dependencies/IncidenceCover.lean}\\
Immutable: \url{https://github.com/alectionary/vibe-math/blob/b7628643bb256bf9c1a319f16f6e19608116cff6/formalization/stopping-bound/classification/dependencies/IncidenceCover.lean}\par}
\smallskip
\medskip
\par\noindent\textbf{\texttt{classification/dependencies/IncidenceDegrees.lean}}\par\nopagebreak
{\footnotesize\noindent Current: \url{https://github.com/alectionary/vibe-math/blob/master/formalization/stopping-bound/classification/dependencies/IncidenceDegrees.lean}\\
Immutable: \url{https://github.com/alectionary/vibe-math/blob/b7628643bb256bf9c1a319f16f6e19608116cff6/formalization/stopping-bound/classification/dependencies/IncidenceDegrees.lean}\par}
\smallskip
\medskip
\par\noindent\textbf{\texttt{classification/dependencies/IncidenceFamily.lean}}\par\nopagebreak
{\footnotesize\noindent Current: \url{https://github.com/alectionary/vibe-math/blob/master/formalization/stopping-bound/classification/dependencies/IncidenceFamily.lean}\\
Immutable: \url{https://github.com/alectionary/vibe-math/blob/b7628643bb256bf9c1a319f16f6e19608116cff6/formalization/stopping-bound/classification/dependencies/IncidenceFamily.lean}\par}
\smallskip
\medskip
\par\noindent\textbf{\texttt{classification/dependencies/IncidenceFamilyCore.lean}}\par\nopagebreak
{\footnotesize\noindent Current: \url{https://github.com/alectionary/vibe-math/blob/master/formalization/stopping-bound/classification/dependencies/IncidenceFamilyCore.lean}\\
Immutable: \url{https://github.com/alectionary/vibe-math/blob/b7628643bb256bf9c1a319f16f6e19608116cff6/formalization/stopping-bound/classification/dependencies/IncidenceFamilyCore.lean}\par}
\smallskip
\medskip
\par\noindent\textbf{\texttt{classification/dependencies/IncidenceFamilySource.lean}}\par\nopagebreak
{\footnotesize\noindent Current: \url{https://github.com/alectionary/vibe-math/blob/master/formalization/stopping-bound/classification/dependencies/IncidenceFamilySource.lean}\\
Immutable: \url{https://github.com/alectionary/vibe-math/blob/b7628643bb256bf9c1a319f16f6e19608116cff6/formalization/stopping-bound/classification/dependencies/IncidenceFamilySource.lean}\par}
\smallskip
\medskip
\par\noindent\textbf{\texttt{classification/dependencies/IncidenceList.lean}}\par\nopagebreak
{\footnotesize\noindent Current: \url{https://github.com/alectionary/vibe-math/blob/master/formalization/stopping-bound/classification/dependencies/IncidenceList.lean}\\
Immutable: \url{https://github.com/alectionary/vibe-math/blob/b7628643bb256bf9c1a319f16f6e19608116cff6/formalization/stopping-bound/classification/dependencies/IncidenceList.lean}\par}
\smallskip
\medskip
\par\noindent\textbf{\texttt{classification/dependencies/IncidenceSource.lean}}\par\nopagebreak
{\footnotesize\noindent Current: \url{https://github.com/alectionary/vibe-math/blob/master/formalization/stopping-bound/classification/dependencies/IncidenceSource.lean}\\
Immutable: \url{https://github.com/alectionary/vibe-math/blob/b7628643bb256bf9c1a319f16f6e19608116cff6/formalization/stopping-bound/classification/dependencies/IncidenceSource.lean}\par}
\smallskip
\medskip
\par\noindent\textbf{\texttt{classification/dependencies/IncidenceSpectral.lean}}\par\nopagebreak
{\footnotesize\noindent Current: \url{https://github.com/alectionary/vibe-math/blob/master/formalization/stopping-bound/classification/dependencies/IncidenceSpectral.lean}\\
Immutable: \url{https://github.com/alectionary/vibe-math/blob/b7628643bb256bf9c1a319f16f6e19608116cff6/formalization/stopping-bound/classification/dependencies/IncidenceSpectral.lean}\par}
\smallskip
\medskip
\par\noindent\textbf{\texttt{classification/dependencies/Spectral.lean}}\par\nopagebreak
{\footnotesize\noindent Current: \url{https://github.com/alectionary/vibe-math/blob/master/formalization/stopping-bound/classification/dependencies/Spectral.lean}\\
Immutable: \url{https://github.com/alectionary/vibe-math/blob/b7628643bb256bf9c1a319f16f6e19608116cff6/formalization/stopping-bound/classification/dependencies/Spectral.lean}\par}
\smallskip
\medskip
\par\noindent\textbf{\texttt{incidence/AuditIncidence.lean}}\par\nopagebreak
{\footnotesize\noindent Current: \url{https://github.com/alectionary/vibe-math/blob/master/formalization/stopping-bound/incidence/AuditIncidence.lean}\\
Immutable: \url{https://github.com/alectionary/vibe-math/blob/b7628643bb256bf9c1a319f16f6e19608116cff6/formalization/stopping-bound/incidence/AuditIncidence.lean}\par}
\smallskip
\medskip
\par\noindent\textbf{\texttt{incidence/Incidence.lean}}\par\nopagebreak
{\footnotesize\noindent Current: \url{https://github.com/alectionary/vibe-math/blob/master/formalization/stopping-bound/incidence/Incidence.lean}\\
Immutable: \url{https://github.com/alectionary/vibe-math/blob/b7628643bb256bf9c1a319f16f6e19608116cff6/formalization/stopping-bound/incidence/Incidence.lean}\par}
\smallskip
\medskip
\par\noindent\textbf{\texttt{incidence/IncidenceCSS.lean}}\par\nopagebreak
{\footnotesize\noindent Current: \url{https://github.com/alectionary/vibe-math/blob/master/formalization/stopping-bound/incidence/IncidenceCSS.lean}\\
Immutable: \url{https://github.com/alectionary/vibe-math/blob/b7628643bb256bf9c1a319f16f6e19608116cff6/formalization/stopping-bound/incidence/IncidenceCSS.lean}\par}
\smallskip
\medskip
\par\noindent\textbf{\texttt{incidence/IncidenceCertificate.lean}}\par\nopagebreak
{\footnotesize\noindent Current: \url{https://github.com/alectionary/vibe-math/blob/master/formalization/stopping-bound/incidence/IncidenceCertificate.lean}\\
Immutable: \url{https://github.com/alectionary/vibe-math/blob/b7628643bb256bf9c1a319f16f6e19608116cff6/formalization/stopping-bound/incidence/IncidenceCertificate.lean}\par}
\smallskip
\medskip
\par\noindent\textbf{\texttt{incidence/IncidenceCore.lean}}\par\nopagebreak
{\footnotesize\noindent Current: \url{https://github.com/alectionary/vibe-math/blob/master/formalization/stopping-bound/incidence/IncidenceCore.lean}\\
Immutable: \url{https://github.com/alectionary/vibe-math/blob/b7628643bb256bf9c1a319f16f6e19608116cff6/formalization/stopping-bound/incidence/IncidenceCore.lean}\par}
\smallskip
\medskip
\par\noindent\textbf{\texttt{incidence/IncidenceCover.lean}}\par\nopagebreak
{\footnotesize\noindent Current: \url{https://github.com/alectionary/vibe-math/blob/master/formalization/stopping-bound/incidence/IncidenceCover.lean}\\
Immutable: \url{https://github.com/alectionary/vibe-math/blob/b7628643bb256bf9c1a319f16f6e19608116cff6/formalization/stopping-bound/incidence/IncidenceCover.lean}\par}
\smallskip
\medskip
\par\noindent\textbf{\texttt{incidence/IncidenceDegrees.lean}}\par\nopagebreak
{\footnotesize\noindent Current: \url{https://github.com/alectionary/vibe-math/blob/master/formalization/stopping-bound/incidence/IncidenceDegrees.lean}\\
Immutable: \url{https://github.com/alectionary/vibe-math/blob/b7628643bb256bf9c1a319f16f6e19608116cff6/formalization/stopping-bound/incidence/IncidenceDegrees.lean}\par}
\smallskip
\medskip
\par\noindent\textbf{\texttt{incidence/IncidenceFamily.lean}}\par\nopagebreak
{\footnotesize\noindent Current: \url{https://github.com/alectionary/vibe-math/blob/master/formalization/stopping-bound/incidence/IncidenceFamily.lean}\\
Immutable: \url{https://github.com/alectionary/vibe-math/blob/b7628643bb256bf9c1a319f16f6e19608116cff6/formalization/stopping-bound/incidence/IncidenceFamily.lean}\par}
\smallskip
\medskip
\par\noindent\textbf{\texttt{incidence/IncidenceFamilyCore.lean}}\par\nopagebreak
{\footnotesize\noindent Current: \url{https://github.com/alectionary/vibe-math/blob/master/formalization/stopping-bound/incidence/IncidenceFamilyCore.lean}\\
Immutable: \url{https://github.com/alectionary/vibe-math/blob/b7628643bb256bf9c1a319f16f6e19608116cff6/formalization/stopping-bound/incidence/IncidenceFamilyCore.lean}\par}
\smallskip
\medskip
\par\noindent\textbf{\texttt{incidence/IncidenceFamilySource.lean}}\par\nopagebreak
{\footnotesize\noindent Current: \url{https://github.com/alectionary/vibe-math/blob/master/formalization/stopping-bound/incidence/IncidenceFamilySource.lean}\\
Immutable: \url{https://github.com/alectionary/vibe-math/blob/b7628643bb256bf9c1a319f16f6e19608116cff6/formalization/stopping-bound/incidence/IncidenceFamilySource.lean}\par}
\smallskip
\medskip
\par\noindent\textbf{\texttt{incidence/IncidenceList.lean}}\par\nopagebreak
{\footnotesize\noindent Current: \url{https://github.com/alectionary/vibe-math/blob/master/formalization/stopping-bound/incidence/IncidenceList.lean}\\
Immutable: \url{https://github.com/alectionary/vibe-math/blob/b7628643bb256bf9c1a319f16f6e19608116cff6/formalization/stopping-bound/incidence/IncidenceList.lean}\par}
\smallskip
\medskip
\par\noindent\textbf{\texttt{incidence/IncidenceSource.lean}}\par\nopagebreak
{\footnotesize\noindent Current: \url{https://github.com/alectionary/vibe-math/blob/master/formalization/stopping-bound/incidence/IncidenceSource.lean}\\
Immutable: \url{https://github.com/alectionary/vibe-math/blob/b7628643bb256bf9c1a319f16f6e19608116cff6/formalization/stopping-bound/incidence/IncidenceSource.lean}\par}
\smallskip
\medskip
\par\noindent\textbf{\texttt{incidence/IncidenceSpectral.lean}}\par\nopagebreak
{\footnotesize\noindent Current: \url{https://github.com/alectionary/vibe-math/blob/master/formalization/stopping-bound/incidence/IncidenceSpectral.lean}\\
Immutable: \url{https://github.com/alectionary/vibe-math/blob/b7628643bb256bf9c1a319f16f6e19608116cff6/formalization/stopping-bound/incidence/IncidenceSpectral.lean}\par}
\smallskip
\medskip
\par\noindent\textbf{\texttt{incidence/dependencies/Spectral.lean}}\par\nopagebreak
{\footnotesize\noindent Current: \url{https://github.com/alectionary/vibe-math/blob/master/formalization/stopping-bound/incidence/dependencies/Spectral.lean}\\
Immutable: \url{https://github.com/alectionary/vibe-math/blob/b7628643bb256bf9c1a319f16f6e19608116cff6/formalization/stopping-bound/incidence/dependencies/Spectral.lean}\par}
\smallskip
\medskip
\par\noindent\textbf{\texttt{source-answer/AuditDenseSourceBridge.lean}}\par\nopagebreak
{\footnotesize\noindent Current: \url{https://github.com/alectionary/vibe-math/blob/master/formalization/stopping-bound/source-answer/AuditDenseSourceBridge.lean}\\
Immutable: \url{https://github.com/alectionary/vibe-math/blob/b7628643bb256bf9c1a319f16f6e19608116cff6/formalization/stopping-bound/source-answer/AuditDenseSourceBridge.lean}\par}
\smallskip
\medskip
\par\noindent\textbf{\texttt{source-answer/AuditSourceAnswer.lean}}\par\nopagebreak
{\footnotesize\noindent Current: \url{https://github.com/alectionary/vibe-math/blob/master/formalization/stopping-bound/source-answer/AuditSourceAnswer.lean}\\
Immutable: \url{https://github.com/alectionary/vibe-math/blob/b7628643bb256bf9c1a319f16f6e19608116cff6/formalization/stopping-bound/source-answer/AuditSourceAnswer.lean}\par}
\smallskip
\medskip
\par\noindent\textbf{\texttt{source-answer/DenseSourceBridge.lean}}\par\nopagebreak
{\footnotesize\noindent Current: \url{https://github.com/alectionary/vibe-math/blob/master/formalization/stopping-bound/source-answer/DenseSourceBridge.lean}\\
Immutable: \url{https://github.com/alectionary/vibe-math/blob/b7628643bb256bf9c1a319f16f6e19608116cff6/formalization/stopping-bound/source-answer/DenseSourceBridge.lean}\par}
\smallskip
\medskip
\par\noindent\textbf{\texttt{source-answer/SourceQuestionAnswer.lean}}\par\nopagebreak
{\footnotesize\noindent Current: \url{https://github.com/alectionary/vibe-math/blob/master/formalization/stopping-bound/source-answer/SourceQuestionAnswer.lean}\\
Immutable: \url{https://github.com/alectionary/vibe-math/blob/b7628643bb256bf9c1a319f16f6e19608116cff6/formalization/stopping-bound/source-answer/SourceQuestionAnswer.lean}\par}
\smallskip
\medskip
\par\noindent\textbf{\texttt{source-answer/dependencies/Spectral.lean}}\par\nopagebreak
{\footnotesize\noindent Current: \url{https://github.com/alectionary/vibe-math/blob/master/formalization/stopping-bound/source-answer/dependencies/Spectral.lean}\\
Immutable: \url{https://github.com/alectionary/vibe-math/blob/b7628643bb256bf9c1a319f16f6e19608116cff6/formalization/stopping-bound/source-answer/dependencies/Spectral.lean}\par}
\smallskip
\medskip
\par\noindent\textbf{\texttt{source-answer/source-ledgers/parent-Audit.lean}}\par\nopagebreak
{\footnotesize\noindent Current: \url{https://github.com/alectionary/vibe-math/blob/master/formalization/stopping-bound/source-answer/source-ledgers/parent-Audit.lean}\\
Immutable: \url{https://github.com/alectionary/vibe-math/blob/b7628643bb256bf9c1a319f16f6e19608116cff6/formalization/stopping-bound/source-answer/source-ledgers/parent-Audit.lean}\par}
\smallskip
\clearpage\section{Complete source chronology and calendar deltas}\label{app:chronology}
This manuscript uses October 24, 2026 in America/Los\_Angeles as its reference date. Civil calendar days to 2026-10-24; exact dates give exact differences. Month/year precision yields inclusive minimum/maximum day ranges; dates are not invented. Date ranges are clipped to the reference date for sources known to exist then. Source age is not unresolved-problem duration.

The central comparison is Table~\ref{tab:chronology}; this ledger also includes inherited sources, software snapshots and held leads. Inspection-only milestones do not become publication dates.

{\small
\begin{longtable}{@{}>{\raggedright\arraybackslash}p{.27\linewidth}>{\raggedright\arraybackslash}p{.47\linewidth}>{\raggedright\arraybackslash}p{.18\linewidth}@{}}
\toprule
\textbf{Source / role} & \textbf{Milestone, date and elapsed calendar days} & \textbf{Access / precision}\\\midrule\endhead
F16: Abiad-Castriota: motivating parameter question; Original question; Prop.5.1 / Remark5.2 / §6 & first and inspected v1: 2026-07-15; 101 days (day) & inspected full text\\[4pt]
F68: Tanner: Minimum Distance Bounds by Graph Analysis; Classical spectral ingredient & earliest inspected manuscript: 2000-07-26; 9586 days (day); journal milestone: 2001; 9063--9427 days (year) & inspected full text\\[4pt]
F67: Zhu-Chugg: Bounds on the Expansion Properties of Tanner Graphs; Classical stopping bound; exact Cor.6 comparison & conference predecessor metadata: 2005-09-28; 7696 days (day); inspected journal version: 2007-12; 6872--6902 days (month) & inspected full text\\[4pt]
F69: Panteleev-Kalachev: Degenerate Quantum LDPC Codes With Good Finite Length Performance; Scalar-factor construction & first v1: 2019-04-04; 2760 days (day); inspected v3: 2021-11-12; 1807 days (day); journal milestone: 2021-11-22; 1797 days (day) & inspected full text\\[4pt]
F70: Leverrier-Zemor: Quantum Tanner codes; Constituent Ramanujan comparison & first v1: 2022-02-28; 1699 days (day); inspected v3: 2022-09-16; 1499 days (day); PDF internal date: 2022-09-19; 1496 days (day) & inspected full text\\[4pt]
F71: Feng-Li: Spectra of Hypergraphs and Applications; Classical source cited by Remark5.2; not a proof premise & publication metadata: 1996-09; 10981--11010 days (month) & abstract only full text gap\\[4pt]
F72: Li-Sole: Spectra of Regular Graphs and Hypergraphs and Orthogonal Polynomials; Classical comparison lead; not a proof premise & publication metadata: 1996; 10889--11254 days (year) & partial web text archival gap\\[4pt]
F73: Prabhu et al.: On the Efficacy of the Peeling Decoder for the Quantum Expander Code; Known stopping-set mechanisms and expansion conditions & first v1: 2025-04-30; 542 days (day); inspected v2: 2025-05-29; 513 days (day) & inspected full text\\[4pt]
F74: Wang-Kulkarni-Poor: Finding All Small Error-Prone Substructures in LDPC Codes; Classical rooted narrowing and branch-and-bound & related first preprint: 2006-05-19; 7463 days (day); inspected current-version date: 2009-04-22; 6394 days (day); journal issue: 2009-05; 6355--6385 days (month) & inspected full text\\[4pt]
F75: Yao-Gokduman-Pfister: Cluster Decomposition for Improved Erasure Decoding of Quantum LDPC Codes; B1 simulations / peeling-cluster comparison & first and inspected v1: 2024-12-11; 682 days (day) & inspected full text\\[4pt]
Connolly2024: Connolly-Londe-Leverrier-Delfosse: Fast erasure decoder for hypergraph product codes; Original stabilizer-support mechanism, credited through inspected F73/F75 & first v1: 2022-08-01; 1545 days (day); v3 metadata: 2024-08-21; 794 days (day); journal publication: 2024-08-27; 788 days (day) & metadata only full text not consulted\\[4pt]
F76: Prabhu-Gayathri-Kumar: On the Analysis of Stopping Sets Associated with the Quantum Hypergraph Product Code; Later comparison lead; no exact theorem claim & conference metadata: 2026; 0--296 days (year) & full text access gap\\[4pt]
RosnesYtrehus2007: Rosnes-Ytrehus: An Algorithm to Find All Small-Size Stopping Sets of Low-Density Parity-Check Matrices; Bibliography-only enumeration lead & conference year: 2007; 6872--7236 days (year) & bibliography only not used as theorem evidence\\[4pt]
Morris2025: Kirsten Morris: Analysis of Graph-based Decoders for Quantum Low Density Parity Check Codes; Discovery trail through existing dissertation inventory & repository source version: 2025-08; 419--449 days (month) & discovery record not new full text screen\\[4pt]
qLDPCOrg/qLDPC: qLDPCOrg/qLDPC pinned snapshot; Software/infrastructure snapshot, not original publication age & inspected commit: 2026-10-23; 1 days (day) & pinned source or infrastructure\\[4pt]
Lcrypto/Stopping-sets-enumerator: Lcrypto/Stopping-sets-enumerator pinned snapshot; Software/infrastructure snapshot, not original publication age & inspected commit: 2023-10-29; 1091 days (day) & pinned source or infrastructure\\[4pt]
leanprover-community/mathlib4: leanprover-community/mathlib4 pinned snapshot; Software/infrastructure snapshot, not original publication age & inspected commit: 2026-10-21; 3 days (day) & pinned source or infrastructure\\[4pt]
X-source-announcement: mathCObot source-paper announcement; Accessible X comparison; not a mathematical answer & post UTC date: 2026-07-16; 100 days (day); authenticated browser local displayed post date: 2026-07-15; 101 days (day) & authenticated post detail inspected\\[4pt]
X-lift-label-announcement: A\_Lectionary lift-label announcement; Distinct mathematical target & post UTC date: 2026-10-24; 0 days (day) & post inspected\\[4pt]
other-code-snapshots: Other inspected GitHub mirrors/README-only code; Secondary discovery/control sources; no theorem or earliest-publication claim & recorded inspection: 2026-10-24; 0 days (day) & inspection date only earliest dates unassessed\\[4pt]
Lean-runtime: Lean 4.36.0-rc1 runtime; Pinned proof checker; release-age claim not made & recorded use: 2026-10-24; 0 days (day) & use date only release date unassessed\\[4pt]
F77: Schwartz-Vardy: stopping distance and stopping redundancy; Classical matrix-dependent stopping convention, codeword mechanism and distinct-column criterion & first preprint metadata: 2005-03-23; 7885 days (day); journal revised manuscript: 2005-12-09; 7624 days (day); inspected journal issue: 2006-03; 7512--7542 days (month) & inspected full text\\[4pt]
F78: Wang-Pryadko: Distance bounds for generalized bicycle codes; Binary GB/CSS construction convention and logical-distance comparison & first and inspected v1 submission: 2022-03-31; 1668 days (day); PDF internal date: 2022-04-01; 1667 days (day) & inspected full text\\[4pt]
F79: Kuo-Ouyang: Degenerate quantum erasure decoding; Quantum-aware decoding comparison; article-in-press main text screened & first preprint metadata: 2024-11-20; 703 days (day); received: 2025-05-31; 511 days (day); accepted: 2026-02-23; 243 days (day); publication / inspected article-in-press: 2026-03-13; 225 days (day); version of record metadata: 2026-04-29; 178 days (day) & inspected full text\\[4pt]
F80: Kashyap-Vardy: Stopping Sets in Codes from Designs; Classical uniform-overlap lower bound, Theorem 1 & reported conference presentation: 2003; 8333--8697 days (year); reported original composition: 2003-09; 8425--8454 days (month); inspected revised author manuscript: 2026-10-24; 0 days (day) & full web text screened direct pdf download403 transcription archived\\[4pt]
F81: Xia-Fu: Minimum Pseudo-Weight and Minimum Pseudo-Codewords of LDPC Codes; Classical overlap/pseudo-weight comparison, Theorems1/3; full17-page primary screen & earliest preprint: 2006-06-12; 7439 days (day); inspected arXivv2 stamp: 2006-11-11; 7287 days (day); PDF internal title-page date: 2019-05-06; 2728 days (day) & inspected full text\\[4pt]
F38: Dehghan-Banihashemi: Cospectral Bipartite Graphs with the Same Degree Sequences but with Different Number of Large Cycles; Earlier cospectral-cycle obstruction; existing23-page source screen reused, no added inventory cohort & first and inspectedv1: 2019-05-31; 2703 days (day); journal volume metadata: 2019; 2489--2853 days (year) & inspected full text\\[4pt]
gideonlsx/bicycle-race: gideonlsx/bicycle-race pinned public repository; README/software comparison; not theorem evidence & inspected commit: 2026-10-11; 13 days (day) & public metadata or readme inspected\\[4pt]
kywukuo/qLDPC\_dec: kywukuo/qLDPC\_dec pinned public repository; README/software comparison; not theorem evidence & inspected commit: 2026-01-26; 271 days (day) & public metadata or readme inspected\\[4pt]
held-lead-1: Generalized Bicycle Codes with Low Connectivity: Minimum Distance Bounds and Hook Errors; Encountered lead, no theorem premise or zero-question claim & inspection only: 2026-10-24; 0 days (day) & metadata or abstract only fulltext unassessed\\[4pt]
held-lead-2: Spectral Theory of Semisimple Bivariate Bicycle Codes; Encountered lead, no theorem premise or zero-question claim & inspection only: 2026-10-24; 0 days (day) & metadata or abstract only fulltext unassessed\\[4pt]
held-lead-3: Logical Operator Decomposition for Distance Analysis of Bivariate Bicycle Codes; Encountered lead, no theorem premise or zero-question claim & inspection only: 2026-10-24; 0 days (day) & metadata or abstract only fulltext unassessed\\[4pt]
held-lead-4: Permutation Decoding and the Stopping Redundancy Hierarchy; Encountered lead, no theorem premise or zero-question claim & inspection only: 2026-10-24; 0 days (day) & metadata or abstract only fulltext unassessed\\[4pt]
held-lead-5: Quantum Maxwell Erasure Decoder for qLDPC codes; Encountered lead, no theorem premise or zero-question claim & inspection only: 2026-10-24; 0 days (day) & metadata or abstract only fulltext unassessed\\[4pt]
held-lead-6: Stopping Sets of Algebraic Geometry Codes; Encountered lead, no theorem premise or zero-question claim & inspection only: 2026-10-24; 0 days (day) & metadata or abstract only fulltext unassessed\\[4pt]
held-lead-7: On the Hardness of Approximating Stopping and Trapping Sets in LDPC Codes; Encountered lead, no theorem premise or zero-question claim & inspection only: 2026-10-24; 0 days (day) & metadata or abstract only fulltext unassessed\\[4pt]
held-lead-8: A class of LDPC codes based on symplectic space over finite fields; Encountered lead, no theorem premise or zero-question claim & inspection only: 2026-10-24; 0 days (day) & metadata or abstract only fulltext unassessed\\[4pt]
F82: Donovan-Rao-Üsküplü-Yaz\i{}c\i{}: QC-LDPC Codes from Difference Matrices and Difference Covering Arrays; Post-lemma primary full-document question screen and exact construction comparison & first and inspectedv1: 2022-05-01; 1637 days (day); journal publication year metadata: 2023; 1028--1392 days (year) & inspected full text\\[4pt]
F83: Chee-Ling-Yin: Optimal Partitioned Cyclic Difference Packings; Post-lemma primary full-document question screen and exact construction comparison & first and inspectedv1: 2010-08-10; 5919 days (day); journal issue metadata: 2010-11; 5807--5836 days (month) & inspected full text\\[4pt]
F84: Gholami-Samadieh: QC-LDPC Codes Based on Finite Set Systems; Post-lemma primary full-document question screen and exact construction comparison & first and inspectedv1: 2014-05-15; 4545 days (day); internal PDF title-page date: 2022-03-28; 1671 days (day) & inspected full text\\[4pt]
F85: Amirzade-Sadeghi: Lower bounds on the lifting degree by difference matrices; Post-lemma primary full-document question screen and exact construction comparison & first and inspectedv1: 2017-09-04; 3337 days (day); internal PDF running draft date: 2018-11-18; 2897 days (day) & inspected full text\\[4pt]
postlemma-held-1: Difference-Spectrum Completion of Multi-Weight Optical Orthogonal Codes; Post-lemma metadata lead, no mathematical premise or zero-question claim & inspection only: 2026-10-24; 0 days (day) & publisher429 and primary PDF403; fulltext unavailable; metadata only no zero question claim\\[4pt]
postlemma-held-2: On the Girth of Tanner QC-LDPC Cycle Codes: An Algebraic Number Theory Approach; Post-lemma metadata lead, no mathematical premise or zero-question claim & inspection only: 2026-10-24; 0 days (day) & publisher429 and primary PDF403; fulltext unavailable; metadata only no zero question claim\\[4pt]
postlemma-held-3: Cyclic balanced incomplete block designs with block size four and arbitrary index; Post-lemma metadata lead, no mathematical premise or zero-question claim & inspection only: 2026-10-24; 0 days (day) & publisher403; abstract metadata only; fulltext unavailable no zero question claim\\[4pt]
postlemma-held-4: Additive, Structural, and Multiplicative Transformations for the Construction of Quasi-Cyclic LDPC Matrices; Post-lemma metadata lead, no mathematical premise or zero-question claim & inspection only: 2026-10-24; 0 days (day) & Secondary search snippet only; primary full version not located/inspected in bounded search; no theorem premise or zero-question claim\\[4pt]
postlemma-held-5: Cyclic Difference Packing and Covering Arrays; Post-lemma metadata lead, no mathematical premise or zero-question claim & inspection only: 2026-10-24; 0 days (day) & Secondary metadata/bibliographic lead only; primary full version not inspected; no theorem premise or zero-question claim\\[4pt]
X-final-search: Final authenticated X source/stopping audit; Read-only bounded query with known-positive control; not absence or theorem evidence & final inspection: 2026-10-24; 0 days (day) & authenticated bounded search read no post sent\\[4pt]
\bottomrule\end{longtable}
}
Inspection-only dates are explicitly distinguished from publication dates. The Kashyap-Vardy author manuscript is undated after its reported 2003 composition; its web transcription hash is not a PDF hash. Supplementary/final-version, citation-index and X access gaps remain named. Date differences are source ages, not durations of unsolved problems.
\clearpage\section{Prior-art coverage and access limitations}\label{app:ledger}
The comparisons use the cited primary theorem statements and their hypotheses.
Public software comparisons include the bicycle-race README and Kuo decoder
repository: their logical-distance and decoder claims do not establish the
stopping-distance guarantees proved here. Accessible X posts include the source
paper announcement and a lift-label result on a different target; neither is
an independent answer to this stopping-bound question.

The recorded GitHub code-search query failed its known-positive control, so its
negative output is not evidence of absence. The focused X search is bounded
and likewise establishes no global absence or historical priority. Named
unavailable full texts, supplementary versions and metadata-only leads are
identified in Appendix~\ref{app:chronology} and the accompanying source ledgers.
Detailed queries, inspected matches and access limits are available in
\href{https://github.com/alectionary/vibe-math/blob/master/formalization/stopping-bound/stronger/prior-art.json}{the prior-art record}
and
\href{https://github.com/alectionary/vibe-math/blob/master/formalization/stopping-bound/stronger/X-GITHUB-FINAL-AUDIT.json}{the public-source audit}.

\end{document}
